% ============================================================
%  R. Nielsen, "Mag. S. Kierkegaards »Johannes Climacus« og
%  Dr. H. Martensens »Christelige Dogmatik«. En undersøgende Anmeldelse."
%  Kjøbenhavn: Forlagt af Universitetsboghandler C. A. Reitzel;
%  Trykt hos Kgl. Hofbogtrykker Bianco Luno, 1849.
%  TRANSCRIPTION (Danish) — from the Google Books scan (Harvard/Widener copy).
%
%  Scan: ~/bibliotek/Nielsen, Rasmus/Kierkegaards_Johannes_Climacus_og.pdf
%        (146 PDF pp., Google Books, clean high-contrast).
%  TYPE: Fraktur (Danish body); Latin/antiqua phrases -> \textit{}.
%        Letterspacing (Sperrsatz) -> \emph{}; image-verify per page by zoom.
%
%  OFFSET (VERIFIED): PDF = printed + 8.
%        printed p.5 = PDF 13; p.6 = PDF 14; ... p.132 = PDF 140.
%        Half-title (cover title) = PDF 7; full title page = PDF 9;
%        text opens on printed p.3 (PDF 11; page number suppressed on the
%        opening leaf, which carries the signature mark "1*" at the foot).
%
%  EXTENT: body printed pp. 3--132 (PDF 11--140). Back matter in the scan:
%        colophon leaf "Trykt hos Kgl. Hofbogtrykker Bianco Luno." (PDF 142).
%        The text closes with the dateline "Kjøbenhavn, d. 18. Sept. 1849."
%        and the signature "R. Nielsen."
%
%  STRUCTURE: a continuous review-essay (no numbered §§, no chapter divisions).
%        It opens by citing the two books under review in display form, states
%        the governing question ("Vil Christendommens Sandhed efter sin Natur
%        være Gjenstand for objectiv Viden?"), and proceeds as unbroken prose
%        with occasional centred rules marking major turns. Rules are
%        reproduced as \begin{center}---\end{center} at the same breaks.
%
%  STATUS: IN PROGRESS. See RESUME-NOTES.md.
% ============================================================
\documentclass[12pt]{book}
\usepackage[utf8]{inputenc}
\usepackage[T1]{fontenc}
\usepackage[danish]{babel}
\usepackage[protrusion=true,expansion=true]{microtype}
\usepackage[margin=1.4in]{geometry}
\usepackage{setspace}
\usepackage{amsmath}
\usepackage{enumitem}
\usepackage{graphicx}    % for \reflectbox (the old Danish »ɔ:« id-est mark)
\DeclareUnicodeCharacter{0254}{\reflectbox{c}}  % ɔ (U+0254) = reversed c
\usepackage{libertinus}
\usepackage{libertinust1math}
\usepackage{textalpha}   % polytonic Greek typed directly, if it occurs later
\usepackage{fancyhdr}
\usepackage[colorlinks=true,linkcolor=black,urlcolor=black,
  pdftitle={Mag. S. Kierkegaards Johannes Climacus og Dr. H. Martensens Christelige Dogmatik},
  pdfauthor={Rasmus Nielsen}]{hyperref}
\setlength{\headheight}{15pt}
\pagestyle{fancy}
\fancyhf{}
\fancyhead[LE,RO]{\thepage}
\fancyhead[RE]{\textit{En undersøgende Anmeldelse}}
\fancyhead[LO]{\textit{Kierkegaards »Johannes Climacus«}}
\setlength{\parindent}{1.5em}
\setlength{\parskip}{0pt}
\onehalfspacing

\begin{document}

% ============================================================
% TITLE PAGE — PDF 9
% ============================================================
\begin{titlepage}
\centering
\vspace*{2.0cm}
{\Large Mag. S. Kierkegaards\par}
\vspace{0.5cm}
{\Huge\bfseries „Johannes Climacus“\par}
\vspace{1.0cm}
{\Large og\par}
\vspace{1.0cm}
{\Large Dr. H. Martensens\par}
\vspace{0.5cm}
{\Huge\bfseries „Christelige Dogmatik.“\par}
\vspace{1.6cm}
\rule{4cm}{0.4pt}\par
\vspace{1.6cm}
{\LARGE\bfseries En undersøgende Anmeldelse\par}
\vspace{0.6cm}
{\large af\par}
\vspace{0.6cm}
{\LARGE\bfseries R. Nielsen,\par}
\vspace{0.25cm}
{\small Professor i Philosophien.\par}
\vfill
{\large\bfseries Kjøbenhavn.\par}
{\small Forlagt af Universitetsboghandler C. A. Reitzel.\par}
{\footnotesize Trykt hos Kgl. Hofbogtrykker Bianco Luno.\par}
{\large\bfseries 1849.\par}
\end{titlepage}

\frontmatter
\mainmatter
\setcounter{page}{3}

% ============================================================
% TEXT — printed pp. 3--132 (PDF 11--140)
% ============================================================

% ---- printed p.3 (PDF 11) ----
Denne Anmeldelse maa begynde med en Undskyldning, at den tillader
sig at sammenstille tvende, i Retning, Aand og Tone saa
ueensartede, Skrifter, som:

\begin{center}
{\large\bfseries Afsluttende uvidenskabelig Efterskrift\par}
{\large til de philosophiske Smuler.\par}
\vspace{0.35em}
{Mimisk-pathetisk-dialektisk Sammenskrift,\par}
{Existentielt Indlæg,\par}
{af Johannes Climacus.\par}
\vspace{0.35em}
{\small Udgiven af\par}
{S. Kierkegaard.\par}
\vspace{0.2em}
{\small Kjøbenhavn 1846, hos C. A. Reitzel.\par}
\vspace{0.9em}
{og\par}
\vspace{0.9em}
{\large\bfseries Den christelige Dogmatik,\par}
{fremstillet af\par}
{\bfseries Dr. H. Martensen.\par}
\vspace{0.2em}
{\small Kjøbenhavn 1849, C. A. Reitzels Forlag.\par}
\end{center}

% ---- printed p.4 (PDF 12) ----
At sammenstille en uvidenskabelig Efterskrift med et
systematisk-videnskabeligt Hovedværk, en experimenterende Humorist
med en christelig Dogmatiker, en flygtig Pseudonym med en
navnkundig Forfatter, vil maaskee hos Mange let opvække en
foruroligende Idee-Association, og gjøre samme uhyggelige Indtryk,
som vilde man nu i hensynsløs Ukritik confundere det Alvorlige med
det Spøgefulde, det Kirkelige med det Theatralske, --- en Prædikant
med en Solodandser. Men, idet denne Mislighed saa stærkt
paatrænger sig, er Anmelderen for Alvor betænkt paa, saavidt det
staaer til ham, at fjerne den, og skal nu strax gjøre Læseren
Regnskab for, i hvilken Henseende de tvende Skrifter vel kunne
antages at belyse hinanden, samt paa hvilken Maade den
Ueensartethed, hvormed de saa brat frastøde hinanden, her lader
sig, --- vel ikke ophæve, men dog i Sammenstillingens Form
betydelig formilde.

Det Spørgsmaal, hvorom Theologien og Philosophien for nærværende
Tid maae siges at bevæge sig, som om det sande og egentlige
Principspørgsmaal, kan kortelig fremsættes i de Ord:

\begin{center}
\emph{Vil Christendommens Sandhed efter sin Natur være Gjenstand
for objectiv Viden?}
\end{center}

At dette Spørgsmaal, naar det opfattes og forstaaes, som det sig
bør, er af gjennemgribende Betydning, er let at indsee. Som
Principspørgsmaal kan det paa ingen Maade holdes hen i undvigende
Ubestemthed,
% ---- printed p.5 (PDF 13) ----
men maa besvares ubetinget, med Ja eller Nei. Besvares
Spørgsmaalet med et ubetinget Ja, da er Skoletheologien i sin gode
Ret, naar den endnu, som hidtil, vedbliver at erklære sig
uundværlig for Kirken. Theologiens Seier over den verdslige Viden
er da et væsentligt Moment i Troens Seier over Vantroen, ligesom
omvendt Theologiens videnskabelige Nederlag ikke blot er et
litterært Uheld for Theologerne, men et smerteligt Tab for den
sande Menighed, en intellectuel Gjenvordighed, der maa opfylde de
Troende med Sorg og Bedrøvelse. Det er da ikke blot paa Dannelsens
og Lærdommens Vegne, men i Troens og Christendommens hellige Navn,
at man, som hidtil, bør holde hart baade ved „den kritiske
Indledning“ og „Hermeneutiken“ og „Apologetiken“ og „den
begrundende, begribende Dogmatik“; men man burde vel endog, især i
disse Tider, conseqvent have en egen Kirkebøn for „den christelige
Videnskab“, en almindelig Bededag for „den kritiske Indledning“,
at dens „Paaskeberegning“ ret maatte styrke vor Tro, og Theologien
holde paa Kirken.

Besvares Spørgsmaalet derimod benegtende, ikke med en undvigende
Vending, som naar der siges, „at Troen er det ene Fornødne, og at
de, der ikke føle Trang til at tænke over deres Tro, gjerne kunne
lade det være“, men med et aabent, ubetinget Nei; indrømmes det
ligefrem, at den christelige Aabenbaring hverken kan eller vil
være Gjenstand for vor objective Viden: da tager Sagen en mærkelig
Forandring; da brister det Led, hvorved Skoletheologien hidtil har
hængt
% ---- printed p.6 (PDF 14) ----
sig paa Kirken; da vil Kirken ogsaa i denne Sag vide at give
Keiseren, hvad Keiserens er, overladende de lærde Discipliner til
deres egen videnskabelige Skjebne.

Her handles altsaa ikke nærmest om en videnskabelig Adskillelse af
den sande og den falske Theologie, men her handles om den
theologiske Videnskabs principielle Forhold til Christendommen;
heller ikke spørges her om Forskjellen imellem den sande og falske
Speculation, men Undersøgelsen gjælder det Spørgsmaal, om
Speculationen overhovedet har nogen religiøs og christelig
Gyldighed, eller om Religionen ikke snarere selv forvandles til
noget ganske Andet, naar den gjøres til Gjenstand for speculativ
Betragtning.

Ved at fastholde det her fremsatte Synspunkt, vil man uden Tvivl
indrømme, at de to anmeldte Skrifter høre nærmere sammen, end det
ved første Øiekast skulde synes; thi medens „Dogmatiken“ gaaer ud
fra den Forudsætning, at Christendommens Indhold bør udlægges,
begrundes og forklares i den objective Videns Former, og altsaa
forsaavidt, skjøndt dog med nogen Usikkerhed, besvarer vort
Spørgsmaal med Ja; er Pseudonymens „uvidenskabelige Efterskrift“
vistnok i ethvert Tilfælde at forstaae som en energisk Protest
imod Nutidens objective Christendomserkjendelse, og besvarer da
Spørgsmaalet med et ubetinget afgjørende Nei.

En betydelig Vanskelighed ligger imidlertid deri, at Pseudonymen
heeltigjennem anvender den indirecte Meddelelse i
dobbeltreflecterede Kategorier, og derved gjør
% ---- printed p.7 (PDF 15) ----
det saa misligt at gjengive hans egen Tankegang i et simpelt
Referat; men deels findes der i Bogen selv ikke faa Partier, hvori
Foredraget nærmer sig til at være „docerende“ (skjøndt hvad der
siges, unegtelig faaer et eget Anstrøg ved den humoristiske
Sammenhæng, hvoraf det udspringer); deels kan jo det Referat, jeg
her agter at give, ganske ligefrem forklares som min Opfattelse,
som det Udbytte, jeg mener at have vundet ved at studere Bogen. Er
det altsaa paa denne Maade forebygget, at min Opfattelse af Bogens
videnskabelige Betydning forvexles% <-- print glyph ambiguous r/x; "forvexles" is the word
{} med Forfatterens derfra maaskee noget afvigende Hovedtanke; da
kan jeg vel ogsaa uden at forvanske Originalens Indhold for min
Deel betjene mig af en ligefrem Meddelelsesform, og derved
muligviis endog opnaae adskillige Fordele. For det Første behøver
jeg da ikke at indlade mig paa Skriftets kunstneriske Anlæg og
eiendommelige Articulation, men kan nøies med at fremhæve de
Punkter, der allernærmest vedkomme Principspørgsmaalet. Dernæst
kan jeg, uden at afvige formeget fra Forfatterens eget Sprog, her
og der modificere Udtrykket saaledes, at det „Pikante“, det
„Urolige“, det „Overgivne“, det „Desultoriske“ ikke faaer Lov at
virke forstyrrende paa den alvorlige Læser.\footnote{Af samme
Grund anseer jeg det heller ikke for hensigtsmæssigt at angive
Paragraph og Pagina ved de enkelte Citater, som naar man ellers
citerer af et Kildeskrift; thi vistnok er denne Efterskrift i sin
Art et Kildeskrift, men just fordi den er et Kildeskrift, der
allerhelst maa læses \textit{in optima forma}, kan det jo
% ---- footnote continues on printed p.8 (PDF 16) ----
bedst overlades Læseren selv at finde de citerede Steder i deres
originale Sammenhæng, hvorhos han da tillige vil kunne have den
Fornøielse at lee og tænke paa een Gang.} Endelig har
% ---- printed p.8 (PDF 16) ----
den hele Sag endnu en Side, man vel kunde ønske lidt nærmere
oplyst. Thi det lader jo næsten, som om den theologiske Opinion nu
vilde vende sig saaledes, at de Indsigelser, den dialektiserende
Pseudonym har reist imod den objective Behandling af Ethiken og
Christendommen, slet ikke skulde fortjene nogen „Litteraturens“
Opmærksomhed, men kun være til for „nogle Enkelte, som, uden at
føle Drift til sammenhængende Tænkning, kunne tilfredsstille sig
selv ved at tænke i Strøtanker og Aphorismer, Indfald og Glimt.“
For da, om muligt, at komme efter, hvorvidt denne Mening er
grundet, vil jeg i al Korthed forsøge, at fremsætte de
paagjældende Hovedpunkter i den Sammenhæng, hvori de have
bekræftet sig for mig. Mislykkes Forsøget, brister Sammenhængen,
bliver Grundtanken veiet og funden for let; da kan Skylden
naturligviis ikke falde tilbage paa en „Efterskrift,“ der anmelder
sig selv som „uvidenskabelig“, men Kritikens Dadel gaaer ud over
mig, der her har troet at see noget Almeengyldigt i Det, som dog,
nærmere beseet, kun er et morsomt Indfald, en lunefuld
Udstrømning af en enkelt vilkaarlig Subjectivitet. Lykkes derimod
Fremstillingen endog blot nogenledes, saa at det dog idetmindste
kan skjønnes, at et Princip er i Bevægelse; da er Fortjenesten
ikke min, men den humoristiske For-
% ---- printed p.9 (PDF 17) ----
fatters, som, efter at være berøvet sin personlige Frihed og sine
glimrende Vaaben, og indtvungen i et docerende Foredrags stive
Form, endnu kan udholde en Dyst med den objective Viden. Udfaldet
blive nu, som det vil, jeg slutter som saa: har Pseudonymen intet
Princip, saa er hans Vittighed ilde anbragt, og „Videnskaben“ bør
da helst sige ham det reent ud, at han ikke tilsidst gaaer hen og
forskruer os den studerende Ungdom. Men er der Noget i, hvad han
siger, har han virkelig sat et Principspørgsmaal i Bevægelse: da
kan man heller ikke vel slippe ham forbi med den Bemærkning, at
han kun er en Individualitet, en genial Særling, der gaaer omkring
som en Undtagelse; thi et Princip er meer end et Individ, og
„Videnskaben“ tør ikke lade et Princip komme bort iblandt
Undtagelserne.

Paa denne Risico maa Anmeldelsen da gaae sin Gang, og Læseren
opfatte dens forskjellige Momenter som Bidrag til en Besvarelse af
det Spørgsmaal:

\begin{center}
{\large\bfseries Har ikke „Johannes Climacus“ dialektisk skjærpet
det christelige Problem?\par}
\end{center}

Skal der være Tale om en christelig-kirkelig Videnskab: da er det
vel indlysende, at denne, som enhver anden Videnskab, maa gaae ud
paa en Løsning af visse bestemte Problemer. Men for at løse et
Problem, maa man først vide, hvor det er, og hvad det er, og
forstaae at fremsætte det med logisk Skarphed. En videnskabelig
Fremstilling, den være forøvrigt saa christelig-kirkeligsindet,
den vil, naar den røber, at Forfatteren
% ---- printed p.10 (PDF 18) ----
enten slet ikke har kjendt Problemet, eller dog kun seet det som i
en svævende Forestilling, er \textit{eo ipso} et svagt Arbeide. I
den theologiske Litteratur, hvor mangehaande lærde Undersøgelser
saa let afdrage Opmærksomheden fra det egentlige Synspunkt, det
sande kritiske Grundpunkt, hvori Troen og Videnskaben mødes, er
det især af Vigtighed at minde herom. I Mathematiken,
Naturvidenskaben og Historien kunne Problemerne maaskee vise sig
paa en mere umiddelbar Maade; men vist er det, at man i Theologien
og Philosophien har omtrent samme Arbeide med at stille Problemet,
som med at løse det. I disse Videnskaber kan da den væsentlige
Opgave stundom blive saaledes omtaaget og fordunklet, at Tænkeren
allerede er mat og udlevet, før han har arbeidet sig igjennem til
Principspørgsmaalet, og opdaget det Vendepunkt, hvorfra den hele
Udvikling nu først egentlig skulde begynde. Vanskeligt har det
hidtil været, endnu vanskeligere bliver det, jo mere de thetiske
Forudsætninger afløses af de dialektiske. I en umiddelbar thetisk
Form kan man endda nogenlunde adskille Philosophien fra
Troeslæren; men, naar Dialektiken saa begynder at gjøre
Grændserne flydende: naar det t.\ Ex.\ hedder, at Forskjellen
imellem det Overnaturlige og Naturlige er flydende, Forskjellen
imellem det Subjective og det Objective flydende, Forskjellen
imellem den explicative og den speculative Begrebsudvikling
flydende, og Forskjellen imellem det Begribelige og det
Ubegribelige ligeledes flydende; da behøves der stor
Skarpsindighed og megen Forsigtig-

% ---- printed p.11 (PDF 19) ----
hed, for at afværge, at alt dette Flydende ikke tilsidst flyder ud
i Eet. Vanskeligheden er her saameget større, som man dog ikke
ligefrem kan vende tilbage til den thetiske Adskillelse, da
Dialektiken unegtelig er Ideebevægelsens høiere Form. Her kræves
da ikke blot en ny Sondring, men en Sondring af høiere Art, en
dialektisk Spænding, hvorunder Problemet igjen kommer tilsyne: her
kræves, at Dialektik corrigeres ved Dialektik. Med denne
Forudsætning vende vi os da nu til Pseudonymen, for at see:

\begin{center}
\emph{a) Hvorledes Problemet fremkommer.}
\end{center}

Uden selv endnu at være kommen ind i Christendommen, kan et
fornuftigt Menneske dog nok forstaae saa Meget, „at den vil gjøre
den Enkelte evig salig, og at den netop hos den Enkelte
forudsætter denne uendelige Interesse for hans Salighed, som
\textit{conditio sine qua non}“. Spørgsmaalet er da ikke nærmest
om Christendommens Sandhed i og for sig, men om „Individets
Forhold til Christendommen.“

Beroer nu Erkjendelsen af Christendommens Sandhed udelukkende paa
Individets sande Forhold til Christendommen, da vil den forfeilede
Bestræbelse tillige være kjendelig paa det usande Forhold ɔ: paa
Misforholdet. Kan det altsaa bevises, at ethvert Forsøg paa at
sikkre sig Christendommens Sandhed ad den objective Vei maa ende
med et uundgaaeligt Misforhold imellem Individet og
Christendommen; da er det jo ogsaa med det Samme beviist, at den
objective Vei fører bort fra
% ---- printed p.12 (PDF 20) ----
Problemet. --- \emph{Den objective Vei deler sig i $\alpha$) den
historiske og $\beta$) den speculative.}

$\alpha$) Betragtes Christendommen som „historisk Actstykke“, og
er det forskende Subject uendelig interesseret i sit Forhold til
Christendommen, saa vil Opgaven jo være at faae en aldeles
paalidelig Efterretning om, hvad der da egentligen er den
christelige Lære. Denne Opgave maa nødvendigviis bringe Subjectet
til Fortvivlelse, da Intet er lettere at indsee, end at den størst
mulige Vished i Forhold til det Historiske kun er en Tilnærmelse,
og at en saadan Tilnærmelse er for Lidet til derpaa at bygge sin
Salighed, ja „saa ueensartet med en evig Salighed, at intet Facit
kan fremkomme.“

Naar saaledes den hellige Skrift betragtes som det sikkre Tilhold,
der afgjør, hvad der er christeligt, og hvad ikke, saa gjælder det
historisk-kritisk at sikkre Skriften. „Man afhandler da her: de
enkelte Skrifters Medhenhøren til Kanon, deres Authentie,
Integritet, Forfatternes Axiopistie, og man sætter en dogmatisk
Garantie: Inspirationen. Naar man tænker paa Engellændernes
Arbeide paa Tunellen, den uhyre Anstrengelse af Kraft, og
hvorledes saa et lille Tilfælde kan i lang Tid forstyrre det Hele
--- saa faaer man en passende Forestilling om dette hele kritiske
Foretagende. Hvilken Tid, hvilken Flid, hvilke herlige Evner,
hvilke udmærkede Kundskaber ere her fra Slægt til Slægt satte i
Reqvisition til dette Underværk. Og dog kan her pludselig en lille
dialektisk Tvivl, der rører Forudsætningerne, for-
% ---- printed p.13 (PDF 21) ----
styrre det Hele, forstyrre den underjordiske Vei til
Christendommen, som man objectivt og videnskabeligt har villet
anlægge, istedetfor at lade Problemet blive til, som det er,
subjectivt.“

At sikkre Bibelen ved Hjælp af Kirken vil ligesaa lidt lykkes, som
omvendt at sikkre Kirken ved Hjælp af Bibelen, naar Sagen holdes
reent objectivt. Foreløbigt er dog dette vundet, at Kirken med det
levende Ord, Troesbekjendelsen og Sakramenterne er et Nærværende,
medens Skriften holder sig i Forbigangenheden. „Kirken“, hedder
det, „bortskjærer al den Bevisen og Bevisen,% <-- print repeats "Bevisen"; zoom-verified, reproduced verbatim
{} som var fornøden i Forhold til Bibelen. At fordre et Beviis af
Kirken for, at den er til, er Nonsens, ligesom at fordre af en
levende Mand Beviis for, at han er til. Kirken er altsaa til. Af
den (som nærværende, som samtidig med Spørgeren, hvorved
Problemet er hævdet Samtidighedens Ligelighed) kan man faae at
vide, hvad der væsentlig er christeligt, thi det er det, som
Kirken bekjender. Rigtigt. Men paa dette Punkt kan Sagen, i
historisk-objectiv Forstand, ikke holde sig. Efterat der nemlig er
sagt om Kirken, at den er til, og at man hos den kan faae at vide,
hvad det Christelige er, udsiges der igjen om denne Kirke, den
nærværende, at den er den apostoliske, at den er den samme, som
har bestaaet i 18 Aarhundreder. Nu kræves der igjen Beviis. Det
eneste Historiske, der er høiere end Beviset, er den samtidige
Tilværelse; enhver Bestemmelse af Forbigangenhed kræver Beviis.
Vil Nogen saaledes sige til En: beviis at Du er til, saa
% ---- printed p.14 (PDF 22) ----
svarer den Anden ganske rigtigt: det er Nonsens. Siger han
derimod: jeg, som nu er til, har været til for over 400 Aar siden,
væsentligen den Samme, saa siger den Første rigtigt: her behøves
et Beviis. Skal der nu føres et historisk Beviis for
Troesbekjendelsens apostoliske Oprindelse og uforandrede Indhold,
saa bliver det igjen en Approximeren, og denne har „den mærkelige
Egenskab, at den kan blive ved, saalænge det skal være.“

Aarhundreders Beviis for Christendommens Sandhed er
utilforladeligt, da det i sit første Ord har forvandlet sig selv
til en Hypothese. „En Hypothese kan blive sandsynligere ved at
holde sig i 3000 Aar, men derfor bliver den aldrig nogen evig
Sandhed, der kan være afgjørende for Ens evige Salighed. Har ikke
Muhamedanismen bestaaet i 1200 Aar? Atten Aarhundreder have ingen
større Beviiskraft end een Dag i Forhold til en evig Salighed;
derimod har atten Aarhundreder og Alt, Alt, Alt det, som kan
fortælles og siges og gjentages i den Anledning, en Adspredelsens
Magt, som distraherer ypperligt.“

Men, kunde En spørge, vil Pseudonymen da maaskee have
Indledningskritiken tilligemed den lærde Philologie reent
afskaffet? --- „Nei, den lærde Philologie er aldeles i sin Ret, og
nærværende Forfatter nærer vist trods Nogen Ærbødighed for, hvad
Videnskab helliger. Men den lærde kritiske Theologie faaer man
derimod ikke noget reent Indtryk af. Hele dens Stræben lider af en
vis bevidst eller ubevidst Dobbelthed. Det seer be-
% ---- printed p.15 (PDF 23) ----
standigt ud, som skulde der af denne Kritik pludselig resultere
Noget for Troen, Noget denne betræffende. Deri ligger
Misligheden“. Hvor vil da Indledningskritiken hen? --- „Har man
Noget bag Øret, saa lad os faae det frem, for i al dialektisk Ro
at tænke det igjennem. Den, der i Hensigt paa Troen forfægter
Bibelen, maa jo have tydeliggjort sig selv, om der af hele hans
Arbeide, hvis det lykkedes efter al mulig Forventning,
resulterede Noget i denne Henseende, paa det han ikke skal blive
siddende i Arbeidets Parenthes og glemme over de lærde
Vanskeligheder det afgjørende dialektiske Claudatur. Den, som
angriber, maa jo ligeledes have gjort sig Rede for, hvis Angrebet
nu lykkedes efter den størst mulige Maalestok, om der saa
resulterede noget Andet end det philologiske Resultat eller i det
Høieste en Seier ved at stride \textit{e concessis}, hvor man, vel
at mærke, kan tabe Alt paa en anden Maade, hvis den gjensidige
Overeenskomst er et Phantom.“

„For at da det Dialektiske kan skee sin Ret, og uforstyrret blot
tænke Tankerne, saa lad os antage først det Ene og saa det
Andet.“

„Altsaa, jeg antager, at det er lykkedes at bevise om Bibelen,
hvad nogensinde nogen lærd Theolog i sit lykkeligste Øieblik har
kunnet ønske at bevise om Bibelen. Disse Bøger høre til Kanon,
ikke andre, de ere authentiske, ere hele, Forfatterne ere
troværdige --- man kan godt sige, det er, som var hvert Bogstav
inspireret (Mere kan man da ikke sige; thi Inspirationen er Troens
Gjen-
% ---- printed p.16 (PDF 24) ----
stand, er qvalitativ dialektisk, ikke til at naae ved en
Qvantiteren). Fremdeles, der er ikke Spor af Modsigelse i de
hellige Bøger. Thi lader os være hypothetisk forsigtige, ymtes der
blot et Ord om noget Saadant, saa er Parenthesen der igjen, og den
philologisk-kritiske Stundesløshed fører En strax paa Afveie ---
Altsaa Alt antaget at være i sin Orden med Hensyn til den hellige
Skrift --- hvad saa? Er saa Den, som ikke havde Troen, kommen et
eneste Skridt nærmere til Troen? Nei, ikke et eneste. Thi Troen
resulterer ikke af en videnskabelig Overveielse, og heller ikke
ligefrem, tvertimod taber man jo i denne Objectivitet den
uendelige personlige i Lidenskab Interesserethed, hvilken er
Troens Betingelse, det \textit{ubique et nusquam}, hvori Troen kan
blive til. Har Den, som havde Troen, vundet Noget i Henseende til
Troens Kraft og Styrke? Nei, ikke det Allermindste, snarere er han
i denne vidtløftige Viden, i denne Vished, der ligger ved Troens
Dør og begjærer efter den, saa farefuldt stillet, at han vil
behøve megen Anstrengelse, megen Frygt og Bæven, for ikke at falde
i Anfægtelse og forvexle Viden med Tro. Medens Troen hidtil i
Uvisheden har havt en gavnlig Tugtemester, vilde den i Visheden
faae sin farligste Fjende. Tages nemlig Lidenskaben bort, saa er
Troen ikke mere til, og Vished og Lidenskab spændes ikke. Hvilken
Lykke da, at denne ønskende Hypothese, den kritiske Theologies
skjønneste Ønske, er en Umulighed, fordi selv den fuldkomneste
Realisation dog kun vil blive en Approximation.
% ---- printed p.17 (PDF 25) ----
Og atter hvilken Lykke for Videnskabsmændene, at Feilen ingenlunde
ligger hos dem! Om alle Engle sloge sig sammen, de kunne dog kun
tilveiebringe en Approximation, fordi i Forhold til en historisk
Viden er en Approximation den eneste Vished --- men ogsaa for Lidt
til derpaa at bygge en evig Salighed.“

„Saa antager jeg det Modsatte, at det er lykkedes Fienderne at
bevise om Skriften, hvad de ønske, saa vist, at det overgaaer den
arrigste Fiendes hidsigste Ønske, --- hvad saa? Har Fienden dermed
afskaffet Christendommen? Ingenlunde. Har han skadet den Troende?
Ingenlunde, ikke det Allermindste. Har han vundet Hævd paa, at
fritage sig selv fra det Ansvar ikke at være en Troende?
Ingenlunde. Fordi nemlig disse Bøger ikke ere af disse Forfattere,
ikke ere authentiske, ikke ere \textit{integri}, ikke ere
inspirerede (dette kan dog ikke modbevises, da det er Troens
Gjenstand), deraf følger jo ikke, at disse Forfattere ikke have
været til, og fremfor Alt ikke, at Christus ikke har været til.
Forsaavidt staaer det den Troende endnu ligesaa frit at antage
det, ligesaa frit, lader os vel agte dette; thi hvis han antog det
i Kraft af noget Beviis, var han isærd med at opgive Troen.
Kommer det nogensinde saavidt, saa vil den Troende altid have
nogen Skyld, dersom han selv har inviteret og begyndt med at
spille Vantroen Seiren i Haanden ved selv at ville bevise.“

$\beta$) \emph{Den speculative Vei.} --- Speculationen bekymrer
sig kun om Christendommen som givet Phænomen,
% ---- printed p.18 (PDF 26) ----
ligegyldig om der gives virkelige enkelte Christne eller ikke, den
antager uden videre, at vi Alle saadan omtrent ere Christne. Der
tales upersonligt: „Speculationen gjør Alt, Speculationen tvivler
om Alt o.\ s.\ v. Den Speculerende taler ikke om sig selv, thi i
Forhold til Speculationen er det egne personlige Selv
ligegyldigt. Men det upersonlige Forhold til Christendommen er et
Misforhold. Er Christendommen en Salighedssag, saa er det jo kun
„to Klasser af Mennesker, der kunne vide Noget om den: De, der
uendeligt interesserede i Lidenskab for deres evige Salighed
troende bygge den paa deres troende Forhold til den; og De, som i
modsat Lidenskab (men i Lidenskab) forkaste den --- de lykkelige
og ulykkelige Elskere. Den objective Ligegyldighed faaer aldeles
Intet at vide.“

Speculationen betragter Christendommen som historisk Phænomen, og
mener at kjende dens Indre af dens Ydre. Det er umuligt. I
Christendommen er det Ydre som Ydre et Ligegyldigt og
Forsvindende. „Tag et Ægtepar; see, deres Ægteskab præger sig
tydeligen ud i det Udvortes; det danner et Phænomen i Tilværelsen
(i det Mindre, ligesom Christendommen verdenshistorisk har
udpræget hele Livet); men deres ægteskabelige Kjærlighed er intet
historisk Phænomen; det Phænomenale er det Ubetydelige, har kun
for Ægtefolkene ved Kjærligheden Betydning; men anderledes seet
(ɔ: objectivt) er det Phænomenale et Bedrag. Saaledes med
Christendommen. Den usynlige Kirke er intet historisk Phænomen;

% ---- printed p.19 (PDF 27) ----
den lader sig slet ikke objectivt betragte som saadan, thi den er
kun i Subjectiviteten.“

For den Speculerende kan nu Spørgsmaalet om hans personlige evige
Salighed slet ikke komme frem, netop fordi hans Opgave bestaaer i
mere og mere at komme bort fra sig selv og blive objectiv, og
saaledes forsvinde for sig selv og blive „Speculationens skuende
Kraft.“ Da Mennesket er en Synthese af det Timelige og det Evige,
vil den Speculationens Salighed, den Speculerende kan have, være
en Illusion, fordi han i Tiden vil være blot evig, vil være
\textit{sub specie æterni}. Heri ligger den Speculerendes
Usandhed. Den i Lidenskab uendelige Interesse for sin personlige
evige Salighed er saaledes høiere end den speculative Salighed,
høiere, fordi den er sandere, fordi den bestemt udtrykker
Synthesen.

„I Forhold til en Iagttagelse, hvor det er af Vigtighed, at
Iagttageren befinder sig i en bestemt Tilstand, gjælder det jo,
at, naar han ikke er i denne Tilstand, saa erkjender han slet ikke
Noget. Han kan nu bedrage En ved at sige, at han er i den
Tilstand, uagtet han ikke er det; men, naar det træffer sig saa
heldigt, at han selv siger, at han ikke er i den fornødne
Tilstand, saa bedrager han Ingen. Dersom nu Christendommen
væsentligen er noget Objectivt, saa gjælder det for Iagttageren at
være objectiv; men, dersom Christendommen væsentligen er
Subjectiviteten, saa er det en Feiltagelse, hvis Iagttageren er
objectiv. I Forhold til al Erkjenden, hvor det gjæl-
% ---- printed p.20 (PDF 28) ----
der, at Erkjendelsens Gjenstand er selve Subjectivitetens
Inderlighed, gjælder det, at den Erkjendende maa være i denne
Tilstand. Men Udtrykket for Subjectivitetens yderste Anstrengelse
er den uendeligt lidenskabelige Interesse for sin evige
Salighed.“

„Den objective Reflexions Vei gjør Subjectet til det Tilfældige,
og derved Existens til et Ligegyldigt, et Forsvindende. Bort fra
Subjectet gaaer Veien til den objective Sandhed, og medens
Subjectet og Subjectiviteten bliver ligegyldigt, bliver den
objective Sandhed det ogsaa, og dette er dens objective Gyldighed,
thi Interessen er, ligesom Afgjørelsen, Subjectiviteten.

„Den objective Reflections Vei fører nu til abstract Tænkning, til
Mathematik, til historisk Viden af forskjellig Art, bestandig
fører den bort fra Subjectet, hvis Tilværelse eller
Ikke-Tilværelse bliver, objectivt, ganske rigtigt, uendelig
ligegyldigt, ganske rigtigt, thi Tilværelse og Ikke-Tilværelse
har, som Hamlet siger, kun subjectiv Betydning. I sit Maximum vil
denne Vei føre til en Modsigelse, og forsaavidt Subjectet ikke
ganske bliver sig selv ligegyldigt, er dette jo kun et Tegn paa,
at hans objective Stræben ikke er objectiv nok; i sit Maximum vil
denne Vei føre til en Modsigelse, at kun Objectiviteten er bleven
til, og Subjectiviteten er gaaet ud, det vil sige: den existerende
Subjectivitet, der har gjort et Forsøg paa at blive, hvad man i
abstract Forstand kalder Subjectiviteten, den abstracte
Objectivitets abstracte Form. Og dog er den Objectivitet, der er
bleven til,
% ---- printed p.21 (PDF 29) ----
subjectivt seet, i sit Maximum enten en Hypothese eller en
Approximation, fordi al evig Afgjørelse netop ligger i
Subjectiviteten.“

\begin{center}
\emph{b) Sandheden, Inderligheden er Subjectiviteten.}
\end{center}

„Man troer i Almindelighed, at det at være subjectiv er ingen
Kunst. Nu, det forstaaer sig, ethvert Menneske er jo ogsaa saadan
et Stykke Subject. Men nu at vorde det, man saadan uden videre er:
ja hvo vilde spilde sin Tid derpaa, det var jo den meest
resignerede af alle Opgaver i Livet. Medens vi Alle ere saadan,
hvad man kalder Subjecter, og arbeide paa at blive objective,
gaaer Poesien bekymret og leder om sin Gjenstand; og dog skal
Poesien netop have Subjectiviteter. Hvorfor tager den saa ikke den
Første den Bedste af vor ærede Midte? Ak nei, han duer ikke, og
vil han ikke gjøre Andet end blive objectiv, saa kommer han aldrig
til at due. Dette synes dog virkelig at tyde paa, at det er en
egen Sag med at være Subject. Hvorfor ere nogle Faa blevne
udødelige som begeistrede Elskere, nogle Faa som høimodige Helte
o.\ s.\ v., dersom Alle i enhver Generation saadan uden videre
vare det ved saadan uden videre at være Subject? Og dog er det at
være Elsker, Helt o.\,s.\,v.\ netop Subjectiviteten forbeholdt,
thi objectivt bliver man det ikke. Poesien vil opløfte os i
Beskuelsen af det Udmærkede --- af hvilket Udmærkede? Ih, af
Subjectiviteten. Saa det er altsaa noget Ud-
% ---- printed p.22 (PDF 30) ----
mærket at være Subjectivitet. Og nu Præsterne! Hvorfor er det kun
et vist Indbegreb af fromme Mænd og Qvinder, til hvis ærværdige
Ihukommelse Talen altid vender tilbage, hvorfor tager Præsten ikke
den Første den Bedste af vor ærede Midte og gjør til Forbillede:
vi ere jo Alle saadan, hvad man kalder Subjecter. Og dog ligger jo
Fromheden netop i Subjectiviteten, objectivt bliver man ikke from.
Præsterne sørge for Sjælene ved at rive os hen i Andagt, medens
Sjælens Længsel eftertragter de Forklarede --- hvilke Forklarede?
Ih, dem, som havde Troen. Men Troen ligger jo i Subjectiviteten,
saa det altsaa er noget Udmærket at være Subjectivitet.“

„Det maa bestandig fastholdes, at det subjective Problem ikke er
Noget om Sagen, men er selve Subjectiviteten. Det gjælder om, at
der slet ikke objectivt bliver Spor af nogen Sag, thi i samme
Øieblik vil Subjectiviteten liste sig fra Noget af Afgjørelsens
Smerte og Krisis ɔ: vil gjøre Problemet lidt objectivt.“

„Den objective Vei mener imidlertid at have en Sikkerhed, som den
subjective Vei ikke har (og det forstaaer sig, Existens, det at
existere og objectiv Sandhed kan ikke tænkes sammen), den mener at
undgaae en Fare, som venter den subjective Vei, og denne Fare er i
sit Maximum Afsindighed. Ved den blot subjective Bestemmelse af
Sandheden, bliver Galskab og Sandhed tilsidst ikke til at
adskille, fordi% <-- print "forbi" (d/b glyph); corrected
{} de begge kunne have Inderligheden. End ikke dette er sandt, thi
Uendelighedens
% ---- printed p.23 (PDF 31) ----
Inderlighed har Afsindighed aldrig; dens fixe Idee er netop et
Slags objectivt Noget, og Afsindighedens Modsigelse netop at ville
omfatte den med Lidenskab. Det, der gjør Udslaget med Hensyn til
Afsindighed, er altsaa atter ikke det Subjective, men den lille
Endelighed, der bliver fixeret, hvilken Uendelighed jo aldrig kan
blive. Eller kan man ikke ogsaa blive gal ved at blive objectiv?
Inderlighedens Udebliven er jo ogsaa Galskab. Den objective
Sandhed som saadan afgjør ingenlunde, at Den, som udtaler den, er
forstandig; tvertimod kan den endogsaa forraade, at Manden er gal,
uagtet det er aldeles sandt, og især objectiv sandt, hvad han
siger.“

Kunde den Existerende virkelig være udenfor sig selv, saa vilde
Sandheden være noget Afsluttet for ham. Men hvor er dette Punkt?
„Jeg-Jeget er et mathematisk Punkt, der slet ikke er til,
forsaavidt kan Enhver gjerne indtage dette Standpunkt, den Ene
staaer den Anden ikke i Veien. Kun momentviis kan det enkelte
Individ existerende være i en Eenhed af Uendelighed og Endelighed,
der er udover det at existere. Dette Moment er Lidenskabens
Øieblik. Den moderne Speculation har opdrevet Alt, for at
Individet objectivt kan komme ud over sig selv, men dette lader
sig slet ikke gjøre. Existensen holder igjen. Den moderne
Speculation lader ringe om Lidenskab; og dog er Lidenskab for en
Existerende netop Existensens Høieste, --- og Existerende ere vi
nu engang. I Lidenskaben er det existerende Subject uendeliggjort
i Phantasiens Evighed, og
% ---- printed p.24 (PDF 32) ----
dog tillige netop allerbestemtest sig selv. Det phantastiske
Jeg-Jeg er ikke Uendelighed og Endelighed i Identitet, thi hverken
det Ene eller det Andet er virkeligt, det er en phantastisk
Overeenskomst i Skyen, en ufrugtbar Omfavnelse, og det enkelte
Jegs Forhold til dette Luftsyn aldrig angivet.“

„Vei-Forskjellen imellem den objective og den subjective
Erkjendelse er nu at udtrykke saaledes:“

„Naar der objectivt spørges om Sandheden, reflecteres der
objectivt paa Sandheden som en Gjenstand, til hvilken den
Erkjendende forholder sig. Der reflecteres ikke paa Forholdet, men
paa at det er Sandheden, det Sande, han forholder sig til. Naar
dette, han forholder sig til, blot er Sandheden, det Sande, saa er
Subjectet i Sandheden.“

„Naar der subjectivt spørges om Sandheden, reflecteres der
subjectivt paa Individets Forhold; naar blot dette Forholds
\emph{Hvorledes} er i Sandhed, saa er Individet i Sandhed, selv om
det saaledes forholdt sig til Usandheden.“

„Lad os tage som Exempel: Guds Erkjendelse. Objectivt reflecteres
der paa, at det er den sande Gud; subjectivt paa, at Individet
forholder sig til et Noget, \emph{saaledes}, at hans Forhold i
Sandhed er et Guds-Forhold. Paa hvilken af Siderne er nu
Sandheden? Mediationen hjælper ikke. En Existerende kan ikke være
to Steder paa engang, kan ikke være Subject-Object. Naar han er
nærmest efter at være to Steder paa en-
% ---- printed p.25 (PDF 33) ----
gang, er han i Lidenskab, men Lidenskab er kun momentviis, og
Lidenskab er netop Subjectivitetens Høieste. Den Existerende, der
vælger den objective Vei, gaaer nu ind i hele den approximerende
Overveielse, der objectivt vil bringe Gud frem, hvilket i al
Evighed ikke naaes, fordi Gud er Subject og derfor kun for
Subjectiviteten i Inderlighed. Den Existerende, der vælger den
subjective Vei, fatter i samme Øieblik den hele dialektiske
Vanskelighed ved, at han skal bruge nogen Tid, maaskee lang Tid,
for at finde Gud objectivt; han fatter denne dialektiske
Vanskelighed i hele dens Smerte, fordi han skal bruge Gud i samme
Øieblik, fordi ethvert Øieblik er spildt, hvori han ikke har Gud.
I samme Øieblik har han Gud ikke i Kraft af nogen objectiv
Overveielse, men i Kraft af Inderlighedens uendelige Lidenskab.
Den Objective generes ikke af saadanne dialektiske Vanskeligheder,
som den: hvad vil det sige, at udvise en heel Forskningens Tid til
at finde Gud.“

„Dersom En, der lever midt i Christendommen, gaaer op i Guds Huus,
i den sande Guds Huus, med den sande Forestilling om Gud i Viden,
og nu beder, men beder i Usandhed, og naar En lever i et afgudisk
Land, men beder med Uendelighedens hele Lidenskab, skjøndt hans
Øie hviler paa en Afguds Billede: hvor er saa meest Sandhed? Den
Ene beder i Sandhed til Gud, skjøndt han tilbeder en Afgud; den
Anden beder i Usandhed til den sande Gud og tilbeder derfor i
Sandhed en Afgud.“
% ---- printed p.26 (PDF 34) ----
„Naar En forsker objectivt efter Udødeligheden, en Anden sætter
Uendelighedens Lidenskab ind paa Uvisheden: hvor er saa meest
Sandhed, og hvo har meest Vished? Den Ene er eengang for Alle
gaaet ind i en Approximeren, som aldrig ender, thi Udødelighedens
Vished ligger jo netop i Subjectiviteten; den Anden er udødelig og
kjæmper netop derfor ved at stride med Uvisheden. Den Ene hjælper
sig ved nogle Beviser, og al hans Interesse falder paa Beviserne,
thi han er jo objectiv, og den objective Viden er ligegyldig mod
det existerende Individ; den Anden sætter hele sit Liv ind paa et
„dersom der er en Udødelighed“, og Uvishedens Smule hjælper ham,
fordi han selv hjælper til med Uendelighedens Lidenskab.“

„Naar Subjectiviteten er Sandheden, maa Sandhedens Bestemmelse
tillige indeholde i sig et Udtryk for Modsætningen til
Objectiviteten, og dette Udtryk angiver da tillige Inderlighedens
Spændstighed. Her er en saadan Definition paa Sandheden:
\emph{den objective Uvished, fastholdt i den meest lidenskabelige
Inderligheds Tilegnelse er Sandheden, den høieste Sandhed, der er
for en Existerende.}“

„Men den givne Bestemmelse af Sandheden er en Omskrivning af Tro.
Uden Risico ingen Tro. Troen er netop Modsigelsen imellem
Inderlighedens uendelige Lidenskab og den objective Uvished.“
% ---- printed p.27 (PDF 35) ----
\begin{center}
\emph{c) For den existerende Subjectivitet er den høieste Sandhed
Paradoxet.}
\end{center}

Paradoxet! Ja, her er Punktet, det kildne Stridspunkt, Problemet,
den gordiske Knude, som „Videnskaben“ ikke kan løse. Maa
Adskilligt af hvad der nylig er sagt om den sande og usande
Gudserkjendelse allerede vække nogen Betænkelighed, blandt Andet,
fordi det er sagt saa dristig og ubetinget, saa maa vel Paastanden
om Paradoxet endnu vække større Anstød og føles uvilkaarlig som et
Stik i Øiet. Jeg tilstaaer, at Intet af hvad jeg ellers har hørt
eller læst om Guds Erkjendelse, har i den Grad tiltrukket og
frastødt mig, som den kjække Paastand, at den høieste Sandhed er
et Paradox. Jeg vedgaaer uden Forbeholdenhed, at Pseudonymen paa
dette Punkt oftere syntes mig at virke langt mere ved et pikant
Indfald end ved en bedre Indsigt, ja han forekom mig endog
undertiden som en bagvendt forheret Tænker, der vilde forvirre
Alt, som et lunefuldt underfundigt Menneske, der ved en
\textit{coup de hasard} vilde snigmyrde „den christelige
Videnskab.“ Men naar jeg saa i et saadant Anfald af Forstemthed
havde kastet Bogen fra mig, og slet ikke vilde have mere med den
Paradoxe at gjøre; vendte han i al Sagtmodighed tilbage igjen, og
begyndte at dialektisere: „Hvorfor afbrød Du Undersøgelsen?“ ---
Den nedbryder Sandheden. --- „Hvilken Sandhed?“ --- Religionens.
Thi den Troende maa dog vel være ligesaa lidenskabelig interesseret
i at vide, at det virkelig er den sande Gud, han tilbeder, som at
han tilbeder Gud i Sandhed. ---
% ---- printed p.28 (PDF 36) ----
„Men kan vel noget Menneske komme til en sand Erkjendelse af Gud
ved at sætte sig i et usandt Forhold til Gud? Forhast Dig ikke;
lad os i Ro tænke Tankerne endnu en Gang! Det sande Forhold til
Gud er jo ikke det objective, hvori Gud bliver ligegyldig for
Individet, idet Individet bliver ligegyldigt for Gud; men det er
det subjective, det er Troens og Tilbedelsens Forhold. Naar nu den
Troende, uendelig bekymret for at faae Vished om den objective
Sandhed (ɔ: Guds Sandhed i sig), ikke tør gaae ud af det
subjective Forhold, for ikke at falde i Usandheden; naar han,
uophørlig arbeidende for at faae det Objective, kun faaer det,
ikke med det Subjective, men som det Subjective: er saa denne hans
Bestræbelse ikke paradox?“ --- Det er unegteligt. --- „Hvorfor
afbrød Du saa Undersøgelsen?“ --- Fordi den stødte min
Subjectivitet. --- „Ja, ved Subjectivitet.“

Det kunde maaskee synes underligt, at jeg ved denne Leilighed
saaledes blotter mig selv og næsten dramatiserer mine subjective
Reflexioner. Hertil kan jeg kun svare, at det ingenlunde skeer,
fordi jeg nu saa gjerne vil stilles offentlig til Skue eller gjøre
mig vigtig ved mine private Bryderier; men --- det er nu engang
gaaet mig saaledes med denne „Johannes Climacus“, og jeg siger det
ligefrem, fordi jeg antager, at det muligviis kunde gaae Andre,
som det er gaaet mig.

Altsaa: Sandheden er Paradoxet. „Naar Subjectiviteten,
Inderligheden er Sandheden, saa er Sand-
% ---- printed p.29 (PDF 37) ----
heden, objectiv bestemmet, Paradoxet, og det at objectivt
Sandheden er Paradoxet, viser netop, at Subjectiviteten er
Sandheden, da jo Objectiviteten støder fra, og Objectivitetens
Frastød eller Udtrykket for Objectivitetens Frastød er
Inderlighedens Spændstighed og Kraftmaaler.“

„Den socratiske Uvidenhed er Udtrykket for den objective Uvished,
den Existerendes Inderlighed er Sandheden. Den socratiske
Uvidenhed er et Analogon til det Absurdes Bestemmelse, kun at der
i det Absurdes Støden fra er endnu mindre objectiv Vished, da der
kun er Vished om, at det er absurdt, og netop derfor uendelig
større Spændstighed i Inderligheden; den socratiske Inderlighed er
et Analogon til Troen, kun at dennes Inderlighed som svarende,
ikke til Uvidenhedens Frastød, men til det Absurdes, er uendelig
dybere. Socratisk er den evige væsentlige Sandhed ingenlunde
paradox% <-- print "parador" (r/x glyph); corrected
{} i sig selv, men kun ved at forholde sig til en Existerende.
Dette er udtrykt i den socratiske Sætning, at al Erkjenden er en
Erindren. Denne Sætning er en Antydning af Speculationens
Begyndelse; men Socrates forfulgte den derfor heller ikke;
væsentligen blev den platonisk.“

„Det uendelig Fortjenstlige ved det Socratiske var netop at
accentuere, at den Erkjendende er Existerende, og det at existere
det Væsentlige. At gaae videre ved ikke at forstaae dette, er kun
en maadelig Fortjeneste. Dette maae vi da altsaa have
\textit{in mente}, og saa see, om Formlen dog ikke lod sig
forandre saaledes, at der virkelig gaaes videre end det
Socratiske.“

% ---- printed p.30 (PDF 38) ----
„Altsaa Subjectiviteten, Inderligheden er Sandheden; gives der et
inderligere Udtryk derfor? Ja, hvis den Tale: Subjectiviteten,
Inderligheden er Sandheden, begynder saaledes: Subjectiviteten er
Usandheden. Man forhaste sig ikke. Speculationen siger ogsaa, at
Subjectiviteten er Usandheden, men siger det lige i den modsatte
Retning, nemlig hen til at Objectiviteten er Sandheden.
Speculationen bestemmer Subjectiviteten negativt hen til
Objectiviteten. Den anden Bestemmelse forhindrer sig selv derimod,
idet den vil begynde, hvilket netop gjør Inderligheden langt
inderligere. Socratisk er Subjectiviteten Usandheden, hvis den
ikke vil fatte, at Subjectiviteten er Sandheden, men f.\ Ex.\ vil
være objectiv. Her derimod er Subjectiviteten, idet den vil
begynde paa at blive Sandheden ved at blive subjectiv, i den
Vanskelighed, at den er Usandheden. Saa gaaer Arbeidet tilbage,
tilbage nemlig i Inderlighed. Det er saa langt fra, at Veien
saaledes skulde være hen til det Objective, at Begyndelsen ligger
endnu dybere i Subjectiviteten.“

„Men Usandheden kan Subjectet ikke være evigt, eller evigt
forudsættes at have været; det maa være blevet det i Tiden eller
blive det i Tiden.“

„Det socratiske Paradox laae i, at den evige Sandhed forholder sig
til en Existerende, men nu har Existensen anden Gang mærket den
Existerende; der er foregaaet en saadan Forandring med ham, at han
ingenlunde kan socratisk erindrende tage sig selv tilbage i Evig-
% ---- printed p.31 (PDF 39) ----
heden. At gjøre dette er at speculere; at kunne gjøre dette, men
at ophæve Muligheden deraf ved at fatte Inderligheden i Existens,
er det Socratiske; men nu er Vanskeligheden den, at det, der
fulgte Socrates, som en ophævet Mulighed, nu er blevet en
Umulighed.“

„Paradoxet fremkommer, naar den evige Sandhed og det at existere
sættes sammen; men for hver Gang det at existere mærkes, bliver
Paradoxet tydeligere og tydeligere. Socratisk seet var den
Erkjendende en Existerende; nu er den Existerende saaledes mærket,
at Existensen har foretaget en væsentlig Forandring med ham.“

„Lad os nu kalde Individets Usandhed \emph{Synd}. Evigt seet kan
han ikke være Synd eller forudsættes evigt at have været i den.
Altsaa ved at blive til (thi Begyndelsen var jo, at
Subjectiviteten er Usandheden) bliver han Synder. Han fødes ikke
som Synder i den Forstand, at han forudsættes som Synder, førend
han fødtes, men han fødes i Synd og som Synder. Dette kunne vi jo
kalde \emph{Arvesynden}. Men har Existensen saaledes faaet Magt
over ham, da er han forhindret i ved Erindring at tage sig selv
tilbage i Evigheden. Var det allerede paradox, at den evige
Sandhed forholder sig til den Existerende; nu er det absolut
paradox, at den forholder sig til en saadan Existerende. Men jo
vanskeligere det er gjort ham erindrende at tage sig ud af
Existens, desto inderligere kan hans Existeren blive i
Existensen.“

„Lad os nu gaae videre, lad os antage, at den
% ---- printed p.32 (PDF 40) ----
evige væsentlige Sandhed selv er Paradoxet. Hvorledes fremkommer
Paradoxet? Ved at den evige væsentlige Sandhed og det at existere
sættes sammen. Naar vi da altsaa sætte det sammen i Sandheden
selv, saa bliver Sandheden et Paradox. Den evige Sandhed er bleven
til i Tiden, dette er Paradoxet. Blev Subjectet i det nærmest
Foregaaende forhindret i at tage sig selv tilbage i Evigheden ved
Synden: nu skal det ikke bekymre sig derover, thi nu er den evige
væsentlige Sandhed ikke bagved, men kommen foran det ved selv at
existere eller have existeret, saa hvis Individet ikke
existerende, i Existensen, faaer fat paa Sandheden, faaer det den
aldrig.“

„Skarpere kan Existens aldrig accentueres, end den nu er bleven
det. Speculationens Svig med at ville erindre sig ud af Existensen
er gjort umulig. Her kan kun være Tale om at fatte dette; enhver
Speculation, der vil være Speculation, viser \textit{eo ipso}, at
den ikke har fattet dette. Individet kan støde alt dette fra sig
og tye til Speculationen; men antage det og saa ville hæve det ved
Speculationen er umuligt, fordi det er lige beregnet paa at
forhindre Speculationen.“

„Naar Socrates troede, at Gud er til, da fastholdt han den
objective Uvished med Inderlighedens hele Lidenskab, og i denne
Modsigelse, i denne Risico er netop Troen. Nu er det anderledes,
istedetfor den objective Uvished, er her Visheden om, at det,
objectivt seet, er det Absurde, og dette Absurde, fastholdt i
Inderlighedens Lidenskab, er Troen. Den socratiske
% ---- printed p.33 (PDF 41) ----
Uvidenhed er som en vittig Spøg i Sammenligning med det Absurdes
Alvor, og den socratiske existerende Inderlighed som en græsk
Sorgløshed i Sammenligning med Troens Anstrengelse.“

„Hvilket er nu det Absurde? Det Absurde er, at den evige Salighed
er bleven til i Tiden, at Gud er bleven til, er født, har voxet
o.\,s.\,v., er bleven til aldeles som et enkelt Menneske, ikke til
at skjelne fra et andet Menneske; thi al umiddelbar Kjendelighed
er for-socratisk Hedenskab og, jødisk seet, Afgudsdyrkelse; og
enhver Bestemmelse af Det, som virkelig gaaer videre end det
Socratiske, maa væsentligen have et Mærke af, at Gud er bleven
til, fordi Tro, \textit{sensu strictissimo}, refererer sig til
Tilblivelse.“

Det Absurde er netop ved det objective Frastød Troens Kraftmaaler
i Inderlighed. Det Absurde er netop Troens Gjenstand og det
Eneste, der lader sig troe. Forsaavidt det Absurde har
Tilblivelsens Moment i sig, vil en Approximations-Vei ogsaa være
den, der forvexler% <-- print "forverler" (r/x glyph); corrected
{} Tilblivelsens absurde Faktum, der er Troens Gjenstand, med et
simpelt historisk Faktum, og altsaa søger historisk Vished for
det, der netop er det Absurde, fordi det indeholder den
Modsigelse, at Det, der kun lige stik imod al menneskelig Forstand
kan blive det Historiske, er blevet det.“

„Christendommen har nu selv forkyndt sig at være den evige
væsentlige Sandhed, der er bleven til i Tiden; den har forkyndt
sig som Paradoxet, og fordret Troens
% ---- printed p.34 (PDF 42) ----
Inderlighed i Forhold til hvad der er Jøder en Forargelse og
Græker en Daarskab --- og Forstanden det Absurde.“

\begin{center}
\emph{d) I Modsætning til Speculationens objective Viden er
Christendommen at betegne som en Existens-Meddelelse.}
\end{center}

„Naar de enkelte Sphærer ikke holdes afgjørende ud fra hinanden,
forvirres Alt. Naar man saaledes er nysgjerrig i Forhold til en
Tænkers Virkelighed, finder det interessant, at man veed Noget
derom o.\ s.\ v., saa er man intellectuelt at dadle, fordi det i
Intellectualitetens Sphære netop er Maximum, at Tænkerens
Virkelighed er aldeles ligegyldig. Men ved saaledes at være
vrøvlevorn i Intellectualitetens Sphære faaer man en forvirrende
Lighed med en Troende. En Troende er netop uendeligt interesseret
i en Andens Virkelighed. Dette er for Troen det Afgjørende, og
denne Interesserethed ikke saadan lidt Nysgjerrighed, men den
absolute Afhængighed af Troens Gjenstand. Problemet udtrykkes
saaledes:

\emph{Individets evige Salighed afgjøres i Tiden ved Forholdet til
noget Historisk, der ydermere er saaledes historisk, at i dets
Sammensætning Det er medoptaget, hvilket ifølge sit Væsen ikke kan
blive historisk, og altsaa maa blive det i Kraft af det Absurde.}

„Troens Gjenstand er ikke en Lære, thi saa er Forholdet
intellectuelt, og det gjælder om ikke at skuffe,
% ---- printed p.35 (PDF 43) ----
men at naae det intellectuelle Forholds Maximum. Troens Gjenstand
er ikke en Lærer, der har en Lære; thi naar en Lærer har en Lære,
er \textit{eo ipso} Læren vigtigere end Læreren, og Forholdet
intellectuelt, hvor det gjælder om ikke at skuffe, men at naae det
intellectuelle Forholds Maximum. Men Troens Gjenstand er Lærerens
Virkelighed, det at Læreren virkelig er til. Troens Svar er derfor
absolut enten Ja eller Nei. Thi Troens Svar er ikke i Forhold til
en Lære, om den er sand eller ikke, ikke i Forhold til en Lærer,
om hans Lære er sand eller ikke, men er Svaret paa Spørgsmaalet om
et Faktum: antager Du, at han virkeligen har været til? Og vel at
mærke, Svaret er med uendelig Lidenskab. --- Troens Gjenstand er
Guds Virkelighed i Betydning af Existens. Men at existere betyder
først og fremmest at være en Enkelt, og derfor er det, at
Tænkningen maa see bort fra Existens, fordi det Enkelte ikke lader
sig tænke, men kun det Almene. Troens Gjenstand er da Guds
Virkelighed i Existens ɔ: som en Enkelt ɔ: at Gud har været til
som et enkelt Menneske.“

„Christendommen er ingen Lære om Eenheden af det Guddommelige og
Menneskelige, om Subject-Objectet, for ikke at nævne de øvrige
logiske Omskrivninger af Christendommen. Dersom nemlig
Christendommen var en Lære, saa var Forholdet til den ikke Troens,
thi til en Lære gives der kun et intellectuelt Forhold.
Christendommen er derfor ingen Lære, men det Faktum, at Gud har
været til som en Enkelt paa Jorden.“
% ---- printed p.36 (PDF 44) ----
„Tro er saaledes ikke en Sinke-Lectie i Intellectualitetens
Sphære, et Asyl for de daarlige Hoveder. Men Tro er en Sphære for
sig selv, og enhver Misforstaaelse strax kjendelig paa, at den
forvandler den til en Lære, og drager den ind i Intellectualitetens
Omfang. Hvad der i Intellectualitetens Sphære gjælder som Maximum,
at blive aldeles ligegyldig mod Lærerens Virkelighed, det gjælder
omvendt i Troens Sphære, dens Maximum er den \textit{quam maxime}
uendelige Interesserethed i Lærerens Virkelighed.“

„Dersom Christendommen var en Lære, vilde den \textit{eo ipso}
ikke danne Modsætning til Speculationen, men være et Moment
indenfor. Christendommen angiver Existens, det at existere, men
Existens og det at existere er netop Modsætningen til Speculation.
Den eleatiske Lære f.\ Ex.\ forholder sig ikke til det at
existere, men til Speculation; der er derfor at anvise den en
Plads indenfor Speculationen.“

„Netop fordi Christendommen ikke er en Lære, derfor gjælder det
den betræffende, at der er en uhyre Forskjel imellem at vide, hvad
Christendom er, og at være Christen.“

„Spørgsmaalet om, hvad Christendom er, maa ikke forvexles% <-- print "forverles" (r/x glyph); corrected
{} med det objective Spørgsmaal om Christendommens Sandhed, som
alt ovenfor er afhandlet. Objectivt lader der sig vel spørge om,
hvad Christendom er, idet den Spørgende vil sætte dette objectivt
hen for sig, og lade det indtil videre staae hen, om den er
% ---- printed p.37 (PDF 45) ----
Sandhed eller ikke. (Sandheden er Subjectiviteten). Den Spørgende
frabeder sig da al velærværdig Travlhed med at bevise dens
Sandhed, samt al speculativ Fremadskynden med at gaae videre; han
ønsker Ro, ønsker hverken Anbefalinger eller Hasten, men at faae
at vide, hvad den er.“

„Eller kan man ikke faae at vide, hvad Christendom er, uden ved
selv at være en Christen? Alle Analogier synes at tale for, at man
kan faae det at vide, og Christendommen selv maa jo betragte dem
som falske Christne, der blot vide, hvad Christendom er. Sagen er
atter her bleven forvirret derved, at man faaer Udseende af at
være en Christen, derved at man er døbt strax som Barn. Men, da
Christendommen kom ind i Verden, eller naar den indføres i et
hedensk Land, da slog den jo ikke og slaaer ikke en Streg over den
samtidige Generation af Ældre og bemægtigede sig de smaa Børn.
Dengang var Forholdet i sin Orden. Da var det vanskeligt, at blive
en Christen, og man befattede sig ikke med at begribe
Christendommen, nu har vi snart naaet Parodien, at det er ingen
Ting at blive Christen; men vanskeligt og en travl Opgave at
begribe den. Herved er alt vendt om, Christendommen forvandlet til
en philosophisk Lære, hvor Vanskeligheden ganske rigtigt ligger i
at begribe, istedenfor at Christendommen væsentligen forholder sig
til Existens, og det at blive Christen er det vanskelige.“

„I Forhold til en Lære er det at forstaae Maxi-

% ---- printed p.38 (PDF 46) ----
mum, og det at blive Tilhænger en underfundig Maade, paa hvilken
Folk, der ikke forstaae Noget, luske sig til at lade, som havde de
forstaaet: i Forhold til en Existens-Meddelelse er det at existere
deri Maximum, og det at ville forstaae (begribe) en underfundig
Udflugt, der vil skulke sig fra Opgaven. At blive Hegelianer er
mistænkeligt, at forstaae Hegel er Maximum; at blive Christen er
Maximum, at ville forstaae (ɔ: begribe) Christendommen er
mistænkeligt.“

„At man kan vide, hvad Christendom er uden at være Christen, maa
da vel bejaes. Noget andet er, om man kan vide, hvad det er at
være Christen, uden at være det, hvilket maa benegtes.“

„Hvad Christendom er, maa da først og fremmest være afgjort, inden
der kan være Tale om nogen Medieren. Herpaa indlader Speculationen
sig ikke, den gaaer ikke saaledes tilværks, at den først
fremsætter, hvad Speculation er, saa hvad Christendom er, for da
at see, om Modsætningerne lade sig mediere; den sikkrer sig ikke
til Vitterlighed de paagjældende stridende Parters respective
Identitet, inden den strider til Forlig. Spørger man den, hvad
Christendom er, svarer den uden videre: den speculative
Opfattelse af Christendommen, uden at bryde sig om, hvorvidt der
er Noget i den Distinction, som distingverer mellem et Noget og
Opfattelsen af Noget, hvilket her synes af Vigtighed for selve
Speculationen; thi, dersom Christendommen selv er Speculationens
Opfattelse, saa bliver der jo ingen
% ---- printed p.39 (PDF 47) ----
Mediation, thi saa bliver der jo ingen Modsætning, og en Medieren
mellem det Identiske er jo dog intetsigende.“

„Men selv naar Speculationen antager en Forskjel mellem
Christendom og Speculation, om ikke af anden Grund, saa for at
have den Fornøielse at mediere, naar den dog ikke bestemt og
afgjørende angiver Forskjellen; saa maa man spørge: er ikke
\emph{Mediationen} Speculationens Idee? Naar der altsaa
\emph{medieres} mellem Modsætningerne, saa er Modsætningerne
(Speculation --- Christendom) ikke lige for det Forligende, men
Christendommen et Moment indenfor Speculationen, og Speculationen
faaer Overvægten, fordi den havde Overvægten, og fordi
Ligevægtens Øieblik, hvor Modsætningerne veiedes mod hinanden,
ikke indtraadte.“

„Dog Misligheden, med at Christendommen skulde være et Moment
indenfor Speculationen, har vel foranlediget Speculationen til at
gaae lidt paa Accord. Speculationen har antaget Titlen:
christelig, har villet anerkjende Christendommen ved at sætte
dette Tillægsord til, ligesom der stundom ved Giftermaal i adelige
Familier dannes et blandet Navn af begge Familiers, eller som naar
Handelshuse forenes i eet Firma, der dog har Begges Navne. Dersom
det nu var saaledes, som man saa let antager, at det at blive
Christen er Ingenting, saa maatte Christendommen rigtignok være
kisteglad, at den havde gjort et saadant godt Parti og var kommen
til Ære og Værdighed næstendeels lige med Speculationen. Dersom
derimod det at blive Christen er den
% ---- printed p.40 (PDF 48) ----
vanskeligste af alle Opgaver, saa synes den Speculerende snarere
at profitere, forsaavidt han med Firmaet vinder det at være
Christen. Men det at blive Christen er virkelig den vanskeligste
af alle Opgaver, fordi Opgaven, skjøndt den samme, varierer sig i
Forhold til det respective Individs Evner. I Forhold f.\ Ex.\ til
det at begribe har et godt Hoved ligefrem Fordelen for den
Indskrænkede, men ikke saaledes i Forhold til at troe. Naar Troen
nemlig fordrer, at han skal opgive sin Forstand, saa bliver det
ligesaa vanskeligt for den Klogeste, som for den meest
Indskrænkede at troe, eller det bliver vel endogsaa vanskeligere
for den Første. Man seer atter Misligheden ved at forvandle
Christendommen til en Lære, hvor det gjælder om at forstaae, thi
derved ligger det at blive Christen i Differensen. Hvad mangler
saa her? Den foreløbige Overeenskomst, hvor hver Parts Status
opgjøres, inden det nye Firma etableres.“

„Blot nu her ikke strax et hurtigt Hoved forklarer en Læseverden,
hvor taabelig min hele Bog er, hvad man mere end tilstrækkeligt
seer deraf, at jeg debiterer noget Saadant, som at Christendommen
ikke er en Lære. Lad os forstaae hinanden. Eet er dog vel en
philosophisk Lære, der vil begribes og speculativt forstaaes, et
Andet en Lære, der vil realiseres i Existens. Naar der i Forhold
til denne sidste Lære skal være Tale om Forstaaelse, saa maa denne
Forstaaelse være at forstaae, at der skal existeres deri, at
forstaae, hvor vanskeligt det
% ---- printed p.41 (PDF 49) ----
er, at existere deri, hvilken uhyre Existens-Opgave denne Lære
sætter den Lærende. Naar det saaledes til en given Tid i Forhold
til en saadan Lære (en Existens-Meddelelse) er blevet almindeligt
at antage, at det at være, hvad Læren byder, er saare let; men det
at forstaae Læren speculativt saa vanskeligt: saa kan En være i
god Forstaaelse med denne Lære (Existens-Meddelelsen), naar han
søger at vise, hvor vanskeligt det er existerende at efterkomme
Læren. I Forhold til en saadan Lære er det derimod en
Misforstaaelse at ville speculere over den.“

„En saadan Lære er Christendommen. At ville speculere over den er
en Misforstaaelse, og eftersom man gaaer videre og videre ad denne
Vei, forskylder man en større og større Misforstaaelse. Naaer man
endeligen det Punkt, at man ikke blot vil speculere, men at man
speculativt har forstaaet den, saa har man naaet Misforstaaelsens
Høieste. Dette Punkt er naaet ved Mediationen af Christendom og
Speculation, og derfor er ganske rigtigt den moderne Speculation
den høieste Misforstaaelse af Christendommen. Naar dette nu
forholder sig saaledes, og naar det fremdeles er givet, at det
nittende Aarhundrede er saa skrækkeligt speculativt, saa er det at
befrygte, at det Ord Lære øieblikkeligt forstaaes om en
philosophisk Lære, der skal og vil begribes. For at undgaae denne
Mislighed har jeg valgt at kalde Christendommen en
Existens-Meddelelse, for ret bestemt at betegne dens
Forskjellighed fra Speculationen.“
% ---- printed p.42 (PDF 50) ----
--- Ved dette Referat bliver Anmeldelsen staaende som et kort
Uddrag af en vidtløftig Forhandling, et lidet Laan fra en stor
Capital. Læseren maa nu selv afgjøre, om han med Anmelderen heri
kan see et Principspørgsmaal bevæge sig taktfast og roligt gjennem
en logisk sammenhængende Tankerække, eller om det Hele kanskee
løser sig op i „Strøtanker og Aphorismer, Indfald og Glimt.“

\begin{center}---\end{center}

Efter Anmeldelsens Ovenstaaende vil det vel uden Vanskelighed
kunne skjønnes, af hvad Grund jeg nu umiddelbart gaaer over til at
omhandle det af Dr.\ H. Martensen nylig udgivne Værk: \emph{den
christelige Dogmatik.}

Da dette dogmatiske Arbeide formeentlig vil være at ansee som et
af de betydeligste i den nyere Theologie, og som et moderne
videnskabeligt Aftryk af den gamle lutherske Læretypus blive
modtaget med Erkjendtlighed i den danske Kirke, ikke blot af de
mange Candidater og yngre Præster, der fornemmelig skylde Dr.\
Martensen deres videnskabelige Dannelse, men ogsaa af Saadanne,
som ellers med nogen kirkelig Sands interessere sig for
Opfattelsen af de christelige Dogmer; kan det med Rette gjøre Krav
paa en fleersidig Opmærksomhed.

Kom det nu blot an paa at bedømme dette omfangsrige Arbeide med
udelukkende Hensyn til Forfat-
% ---- printed p.43 (PDF 51) ----
terens Talent, Studium og Fremstillingsgave, da vilde jeg for mit
Vedkommende kun have Anledning til paa Ny at anerkjende den Omsigt
i Stoffets Behandling, den simple Klarhed i Formen og Reenhed i
Udtrykket, der i dette Skrift, som i „Mester Eckart“, tiltaler
Læseren fra Begyndelsen til Enden. Her er jo vistnok i denne
christelige Dogmatik, seet fra Formsiden, ogsaa noget Behageligt,
noget Tiltalende, Noget, der i visse Stemninger virker, om just
ikke strengt overbevisende, saa dog mildt beroligende. Her er i
denne christelige Dogmatik noget Fleersidigt, noget Forsonende,
noget Fredsommeligt, noget Føieligt. Her er i denne christelige
Dogmatik, naar den betragtes som systematisk Heelhed, tillige
noget vist Kunstnerisk, noget Architektonisk, Noget, der paa een
Gang ligesom formæler Universitetsbygningen med Slotscapellet,
Videnskaben med Religionen: skarp Forstand og naiv Enfoldighed,
moderne Reflexion og mystisk Umiddelbarhed, Begrebets Klarhed og
Mysteriets Dunkelhed: Alt smelter sammen i Stilens Eenhed, i
festlig Ro, høitidelig Dæmring, som Dagens Skjær med Skinnet af et
Alterlys. Det er et ædelt Talent, en prægnant Naturgave, hvorved
Dr.\ Martensen, som Universitetslærer, har beaandet saa mange
Disciple, og, som geistlig Taler, opbygget saa mange Tilhørere.

Som sagt, kom det blot an paa at fremhæve og anerkjende den
skjønne Formside, da vilde jeg i det her omhandlede Skrift finde
Steder nok, der kunde tjene mig som Bilag. Heller ikke vil Nogen,
der har lagt Mærke
% ---- printed p.44 (PDF 52) ----
til mine videnskabelige Forsøg, let være uvidende om, hvor megen
Priis jeg i saa Henseende sætter paa Dr.\ Martensens Talent og
Smag. Men det Spørgsmaal, der i denne Anmeldelse skal forhandles,
er nu intet æsthetisk Formspørgsmaal, men et afgjørende
Principspørgsmaal, og jeg vender derfor tilbage til mit Første:
\emph{Vil Christendommen efter sin Natur være Gjenstand for
objectiv Viden?} hvilket Spørgsmaal i denne Sammenhæng kan
fastholdes med den Begrændsning: \emph{„Vil Christendommen være
Gjenstand for dogmatiserende Speculation?“} --- Dette Spørgsmaal,
der nu er blevet mig til et Slags litterairt
Samvittighedsspørgsmaal, vil jeg da i denne Anmeldelse ikke mere
tabe af Syne.

Da Professor Martensen optraadte ved Universitetet, og vakte hos
den studerende Ungdom den glade Forventning, at Christendommen nu
vilde lade sig forklare i Speculationen, havde han allerede (1837)
udgivet sit første theologiske Skrift: „Om den menneskelige
Selvbevidstheds Autonomie i den nyere Tids dogmatiske
Theologie“,\footnote{Oversat af L. B. Petersen.} og i dette Skrift
angivet Forudsætningen for en christelig Speculation. Under samme
Forudsætning udarbeidede jeg da ogsaa nogen Tid efter (1840) min
første Afhandling: „Den speculative Methodes Anvendelse paa den
hellige Historie“\footnote{Oversat af Cand.\ theol.\ B. C.
Bøggild.}, et Ungdomsarbeide, hvorved
% ---- printed p.45 (PDF 53) ----
jeg agtede at forberede en tilsvarende Fremstilling af Jesu Liv.
En Fremstilling af Jesu Liv har jeg da nu ogsaa
udarbeidet\footnote{Evangelietroen og den moderne Bevidsthed,
Forelæsninger over Jesu Liv. 1 Deel, 1849, hos C. A. Reitzel.};
men Udførelsen svarer ikke til Forberedelsen: Standpunktet er
blevet et andet. Ingenlunde, at jeg har opgivet Speculationen,
endnu mindre Christendommen, men jeg har, altid forsmaaende
Middelveislærens Modificationer, opgivet min tidligere Forventning
om begges Sammensmeltning, og altsaa i dette Punkt forandret
Anskuelse, eller, som det hedder, skiftet Standpunkt.

Anderledes er det gaaet Dr.\ Martensen, der, som organisk Genius,
er bleven sin første Kjærlighed tro og ikke har skiftet
Standpunkt: hans Skrift om „Autonomien“ svarer nu til hans
Dogmatik, som Forberedelsen til Udførelsen, som Knoppen til
Blomsten, Prophetien til Opfyldelsen.\footnote{„Hvorledes jeg har
ladet mig belære af Tidernes Tegn i Henseende til det saa meget
omtvistede Spørgsmaal om Erkjendelsens Forhold til Troen, maa
Gjerningen udvise. At tale udførligere herom i et Forord, anseer
jeg for unyttigt, da jeg har gjort det i Indledningen, og Alt
tilsidst dog kommer an paa Udførelsen. Kun maa det være mig
tilladt at sige, at jeg i den Vending, de philosophiske Bevægelser
have taget i det sidst forløbne Tidsrum, ikke har fundet nogen
Anledning til at opgive den Tanke, som ligger til Grund for mit
første theologiske Skrift: „Om den menneskelige Selvbevidstheds
Autonomie i den nyere Tids dogmatiske Theologie“, den Tanke, at
det ingenlunde er Religionen, der skal laane sin Vægt og
% ---- footnote continues on printed p.46 (PDF 54) ----
Betydning af den speculative Tænkning, men at det er den
speculative Tænkning, der trænger til Religionen, til Guds
Aabenbaring som til sit Princip.“ Hin Selvbevidsthedens Autonomie,
som jeg dengang (1837) søgte at paapege, er efter den Tid med
forbausende Hurtighed gaaet over i nye
Udviklingsformer.“\ldots{} Af Forordet til d.\ chr.\ Dogmatik.}
Vel har Dr.\ M. i adskillige Henseen-
% ---- printed p.46 (PDF 54) ----
der ladet sig belære af „Tidernes Tegn“; men, hvorledes det med
denne Laden sig belære nærmere hænger sammen, er foreløbig antydet
i Forordet til Dogmatiken, og af dette Forord hensætter jeg da her
Følgende:

„At en i sig selv sammenhængende theologisk Tænkning, ja
theologisk Speculation, baade er mulig og nødvendig, er en
Overbeviisning, jeg haaber at kunne fastholde, selv om det skulde
kunne eftervises, at mit eget Arbeide var forfeilet. Ligesaalidt
som en negativ Speculation har kunnet bringe mig til at opgive
den, ligesaalidt have de Indsigelser kunnet det, som jeg i den
senere Tid har hørt fremsættes i Troens Navn. Psychologisk
betragtet maa det vistnok findes aldeles forklarligt, at der
nuomstunder ere Mange, som have Mistillid til al Troes-Videnskab
og mene, at Videnskaben alene hører Fritænkeriet og Hedenskabet
til. Til enhver Tid gjælder det jo og, at ikke Videnskaben, ikke
Theologien, men kun Troen er det ene Fornødne for os Alle. De, som
altsaa ikke føle nogen Trang til at tænke over deres Tro, kunne
være i deres gode Ret, naar de for deres eget Vedkommende betragte
Erkjendelsen som ufornøden. Og de, som ikke føle Drift til
sammenhængende Tænk-

% ---- printed p.47 (PDF 55) ----
ning, men kunne tilfredsstille sig selv ved at tænke i Strøtanker
og Aphorismer, Indfald og Glimt, kunne ogsaa være i deres Ret,
naar de for deres eget Vedkommende betragte sammenhængende
Erkjendelse som ufornøden. Men, naar det i den senere Tid begynder
at udpræges til et Slags Dogma, at den Troende aldeles ikke kan
have nogen Interesse af at søge sammenhængende Erkjendelse af Det,
som er ham det Høieste; at den Troende ikke kan ville indlade sig
paa nogen Speculation, fordi al Speculation kun er kosmisk, det
vil sige: verdslig og hedensk; at den Troende maa betragte
Begrebet om Troes-Videnskab som en Selvmodsigelse, der ophæver den
ægte Christendom o.\ s.\ v., da tilstaaer jeg, at saadanne
Sætninger, selv naar jeg har hørt eller seet dem fremsatte med det
Aandriges Paradoxie, for mig have havt saa lidet Overbevisende, at
jeg deri kun kan see en stor Misforstaaelse og en ny, eller
rettere sagt en gammel Misviisning. Saa Meget har jeg deraf kunnet
lære: at ligesom der var en Tid, da vi havde ikke Faa iblandt os,
som altfor tidligt vare blevne absolute i Viden, saaledes fattes
det nu igjen ikke paa Saadanne, som altfor tidligt ere blevne
absolute i Troen.“

Saaledes har Dr.\ Martensen altsaa ladet sig belære af „Tidernes
Tegn“, saa let er han bleven færdig med de Indsigelser, der i
Troens Navn ere reiste mod Speculationen, paa denne Maade tager
han Ordet for den sammenhængende Tænkning og vidner for den
objective Viden.
% ---- printed p.48 (PDF 56) ----
Det er virkelig meget smukt af Dr.\ Martensen, at han endnu vil
vidne for den „christelige Speculation“ endog paa en Tid, da saa
Mange ere frafaldne, thi dette taler jo for, at han selv maa have
en inderlig fast og velbegrundet Overbeviisning om Speculationens
Nødvendighed; det er ligeledes smukt, at Dogmatikeren, der lever i
en Tid, hvor de urolige Tanker kun altfor let adsprede sig i
„Aphorismer, Indfald og Glimt“, alvorlig minder os om, hvor
nødvendigt det er at søge en sammenhængende Erkjendelse af Det,
som bør være Mennesket det Høieste. Vi skylde Dr.\ Martensen vor
Tak saavel for Vidnesbyrdet som for Paamindelsen; men hvorledes
skulde vi vel takke ham paa en værdigere Maade, end naar vi efter
ringe Evne stræbe at sætte os ind i hans dybsindige Speculation,
og med al Flid randsage, hvorledes det da nu staaer sig med hans
egen „sammenhængende Erkjendelse.“

Men --- Opgaven er vanskelig og Løsningen tvivlsom: tomme
Paastande og stive Forsikkringer bevise jo Intet; det er da vel
det Bedste, at lade Undersøgelsen svæve som en Sag, der svæver for
Retten, og imidlertid holde sig nær til det Spørgsmaal:

\begin{center}
{\large\bfseries Har ikke „den christelige Dogmatik“ udialektisk
omgaaet det christelige Problem?\par}
\end{center}

Hvorledes man end iøvrigt vil bestemme Christendommens og
Speculationens gjensidige Forhold: dette maa man dog vel
indrømme, at de hver i sin Form idet-
% ---- printed p.49 (PDF 57) ----
mindste gjøre Fordring paa at udtrykke det Absolute.
Christendommen meddeler det Absolute i Form af Tro, Speculationen
i Form af Viden. Den, der holder sig til Speculationen, vil altsaa
igjennem Viden staae i et absolut Forhold til det Absolute, og
forsaavidt være „absolut i Viden.“ Den derimod, der holder sig til
Christendommen, vil igjennem Troen staae i et absolut Forhold til
det Absolute, og maa da vel forsaavidt være „absolut i Troen.“

At erklære sig „absolut i Viden“ er en Anmasselse, forsaavidt den
Speculerende dermed vil sige sig at have en i alle Henseender
absolut og fuldkommen Viden af det Absolute; men forsaavidt han
dermed vil udtrykke, at det Absolute er og maa være ham det eneste
og eneherskende Princip i al hans Tænken og Viden, er dette saa
langt fra at være nogen Anmasselse, at det meget mere er en
pligtskyldig og nødvendig Selvbegrændsning. At være „absolut i
Viden“ er da det Samme som at være streng og samvittighedsfuld og
nøieregnende med sin Viden af det Absolute. Den Vidende tør ikke
vente nogen Bistand andetsteds fra, t.\ Ex.\ fra Troen, saa at det
Absolute skulde bekræfte sig for ham halvt i Troen og halvt i
Viden; thi i saa Fald har han krænket sit Princip og forqvaklet
sin Videnskab, ligesom en Mathematiker jo har krænket sit Princip
og forqvaklet sin Videnskab, naar han istedetfor stringente
Beviser argumenterer af den umiddelbare Erfaring. Om end den
Vidende, som existerende Individ, kun trænge til mange andre
% ---- printed p.50 (PDF 58) ----
Oplysninger og en ganske anden Opfattelse af det Sande, end den,
der kan gives ham i Videns absolute Form, saa vedkommer alt Sligt
dog ikke hans Viden. Som Vidende har han slet ingen anden
Interesse og ingen anden Trang end den at faae det Absolute at
vide. Kun paa disse Vilkaar kan Speculationen udvikle sig og hævde
sin Ret og gjøre Fremskridt; kun paa disse strenge Vilkaar have de
store Tænkere, en Fichte, en Schelling, en Hegel arbeidet sig
frem, og ere, hver efter sit Standpunkt, blevne „absolute i
Viden.“

At ville gaae og gjælde for „absolut i Troen“ er en Anmasselse,
hvis den Enkelte dermed vil udgive sig for at kunne fuldkomme al
Troens Gjerning, at have en i alle Maader absolut og ubetinget
Tro, den Tro, som flytter Bjerge; men, naar det derimod forstaaes
saaledes, at det er det Samme, som at gjøre Troen til absolut
ufravigelig Betingelse for den existentielle Tilegnelse af
Christendommens Absolute; da er dette saa langt fra at være en
Anmasselse, at det endog ligefrem er den Troendes Pligt og
Skyldighed. Er Troen den eneste rette Form og Betingelse,
hvorunder Christendommens absolute Sandhed kan tilegnes: saa maa
jo ogsaa Troesprincipet i sig selv være et absolut Princip, et
Princip, der kun krænkes og forqvakles, naar det skal
understøttes af et Andet, t.\ Ex.\ af den speculative Viden. Kan
end den Troende, som tænkende Aand (og da i Særdeleshed som
virkelig Genius), føle Trang til megen anden Kundskab og en ganske
anden Erkjendelse,
% ---- printed p.51 (PDF 59) ----
end den, der ligefrem er indeholdt i Troens Absolute: saa staaer
dette dog fast, at han som Troende hverken har eller kan have
anden Trang eller anden Interesse eller anden Opgave, end den, at
arbeide for Tilegnelsen og fuldkomme Troens Gjerning. Kun paa
disse Vilkaar kan Troeserkjendelsen klares og fremmes og hævde sin
Ret; kun paa disse strenge Vilkaar have de store Troesvidner
ligefra Paulus indtil Luther lidt og stridt, og ere, hver efter
sin Aandens Gave, blevne „absolute i Troen.“

Ingen mene, at jeg med disse Bemærkninger indbilder mig at sige
noget Nyt, eller dog Noget, som Dr.\ Martensen ikke veed.
Langtfra! Dr.\ M. veed det Altsammen ligesaa godt som jeg, og kan,
naar han vil, desuden sige det langt smukkere. Idetmindste har han
tidligere vidst, at tale saa absolut om det Absolute, at det var
en sand Glæde at høre paa. Men hvad jeg her siger, det siger jeg i
Tvivlraadighed, fordi det forekommer mig betænkeligt, fordi det
forekommer mig, at min høiærværdige Ven, „belært af Tidernes
Tegn,“ har i en vigtig Sag begyndt at forandre sit Sprog, uden dog
at have forandret sit Standpunkt. See, det er en Omstændighed, der
gjør mig tvivlraadig, saa jeg næsten ikke veed, hvad jeg skal
sige, en betænkelig Vending, der gjør mig Talen saa forblommet, at
jeg endog kunde fristes til at adspørge „den kritiske Indledning“
og søge Raad hos „Hermeneutiken.“
% ---- printed p.52 (PDF 60) ----
Hvormeget er det da, Dr.\ M. har kunnet lære af „Tidernes Tegn“?
Svar: „at ligesom der var en Tid, da vi havde ikke Faa iblandt os,
som altfor tidligt vare blevne absolute i Viden, saaledes fattes
det nu igjen ikke paa Saadanne, som altfor tidligt ere blevne
absolute i Troen.“ Altsaa: Saameget! idetmindste er det Saameget,
der siges med Bestemthed; thi hvad der i det Følgende siges, er
ikke nær saa bestemt. Men hvormeget er nu dette „Saameget“, naar
vi see nærmere til? Ja det er just det, jeg ikke kan udfinde.
Tager jeg de tvende Udtryk: „absolut i Viden“ og „absolut i
Troen“ saaledes, at de sigte paa den Opblæsthed, der kun altfor
gjerne ledsager saavel de speculative Opsving, som de religiøse
Opvækkelser, saa bliver Meningen denne: „da Speculationen
herskede, havde vi ikke Faa iblandt os, som indbildte sig at være
speculativt alvidende;“ og hvis det skal forstaaes paa denne
Maade, saa har Dr.\ Martensen sandelig Ret i, at slige Personer
vistnok maae være komne „altfor tidligt“ med deres Viden. Men
Fortolkningen passer ikke; thi hvormegen Umodenhed der end i sin
Tid kan have yttret sig hos den speculerende Ungdom: saa har det
dog vist kun været meget Faa (ifald der ellers har været Nogen),
som bildte sig ind at have naaet den speculative Alvidenhed. Vi
skylde dog vel Dr.\ M. saamegen Agtelse, at vi ikke gaae hen og
fortolke hans Ord paa en saa ufornuftig Maade, at Alt, hvad han
har lært af „Tidernes Tegn“, opløser sig i det rene Nonsens. Dette
kan desuden saameget lettere afvær-
% ---- printed p.53 (PDF 61) ----
ges, som den her foreslaaede Fortolkning, der saa maadeligt passer
paa det første Led (de absolute i Viden), til al Lykke slet ikke
lader sig anvende paa det andet (de absolute i Troen); thi de, der
paa denne Maade skulde være komne „for tidligt“, maatte jo være
rene religiøse Sværmere; men af Saadanne findes der dog vel endnu
ingen iblandt os.

Men naar jeg nu forstaaer Udtrykkene: „absolut i Viden“ og
„absolut i Troen“ om en absolut Anerkjendelse saavel af Videns som
af Troens Princip, er jeg endda lige nær. Hvorledes kan Dr.\ M.
betragte det som en ny og dog gammel Misviisning, at Nogle ville
være absolute i Videns Princip, Andre absolute i Troens, da Dr.\
M. jo selv, ifølge sit Standpunkt, har saa godt som forpligtet sig
til at være absolut i dem begge?

At lægge det hele Eftertryk paa „altfor tidligt“, nytter ikke; thi
at komme for tidligt til en absolut Anerkjendelse af et absolut
Princip vilde jo være det samme som at komme for tidligt til
Erkjendelse af en væsentlig Sandhed. Men saa er det dog noget
forunderligt at klage over, at Menneskene komme for tidligt til
Sandheds Erkjendelse; man pleier ellers at klage over, at de komme
for sildigt. At Dr.\ M. maa have lært noget Vist af „Tidernes
Tegn“, ja noget ganske Bestemt, er jeg overbeviist om, da han selv
siger det; men hvad det nu egentlig er for Noget, og hvad dette
Noget egentlig skal betyde, er jeg ikke istand til at udfinde,
medmindre det skulde kunne opdages i det Følgende.
% ---- printed p.54 (PDF 62) ----
„Saavidt jeg nemlig“ --- vedbliver Dr.\ M. --- „formaaer at
skjønne, er der kun Een, som fuldkomment svarer til Begrebet om
den Troende, nemlig hele den almindelige Kirke. Enhver af os
Enkelte besidder kun Troen i et vist begrændset Maal, og maa vel
vogte os for, at vi ikke gjøre vort eget individuelle, maaskee
noget eensidige, maaskee endogsaa noget sygelige Troesliv til
Regel for alle de Troende. Kun hiin store Enkelte, der, som
Apostelen lærer os, gjennem Tiderne skal opvoxe til en fuldkommen
Mand, kan fuldstændigt virkeliggjøre Begrebet om den Troende i
dets hele Aandelighed, dets hele Alsidighed og rige
Mangfoldighed, kan ene besidde Troens og alle Naadegavers Fylde.
Men om hiin store Enkelte vide vi efter Apostelens Vidnesbyrd, at
han er i en fortsat Fremskriden henimod Troens og Erkjendelsens
Eenhed. (Eph.\ 4, 13).“

Dersom det nu ikke var Dr.\ Martensen, der talte til os i disse
mærkværdige Linier, skulde jeg for min Part snart blive færdig med
Forklaringen; thi dybe Seere, der tolke „Tidernes Tegn“, uden at
man derfor kan sige, hvad det er, de egentlig have seet; kloge
Folk, som i principielle Stridigheder give gode Raad, uden at man
just deraf tør slutte, at de egentlig vide, hvorom der strides;
litteraire Autoriteter, der i fornem Overlegenhed advare Samfundet
mod visse Enkeltes Udskeielser, uden at de derfor finde det
passende at nævne disse Enkelte ved Navn: Sligt er endda hverken i
Theologien eller Philosophien noget saa ganske Usædvanligt, og man

% ---- printed p.55 (PDF 63) ----
behøver ikke at tage sig Sagen videre nær. Men, naar en Mand, der
vil Noget, og veed, hvad han vil, en Mand, der veier sine Ord, og
veed, hvad han siger, en Mand som Forfatteren af „den christelige
Dogmatik“, der baade kan og vil klare os et theologisk
Grundspørgsmaal, begynder med at tale i Gaader, da bliver man
uhyggelig tilmode, usikker paa sig selv og geraader i en piinlig
Forlegenhed. Skulde her iblandt os virkelig findes nogle Enkelte,
som af en Feiltagelse have forvexlet sig selv med den „hellige,
almindelige Kirke“, nogle Enkelte, der bilde sig ind at have den
fuldkomne Tro; nogle selvraadige Enkelte, der ikke blot ville
løsrive sig fra Kirken, men ovenikjøbet gjøre deres eget
„individuelle, eensidige og sygelige Troesliv til Regel for alle
de Troende“? Men hvordan hænger Dette dog sammen? Hvilke ere disse
Enkelte? hvor befinde de sig? hvad kalder man dem? --- Jeg seer og
seer --- paa „Tidernes Tegn“; jeg seer en talløs Vrimmel af
forskjellige, i Sind og Tanke høist forskjellige Enkelte, men
disse bestemte Enkelte kan jeg ikke see. „Hiin store Enkelte“, om
hvilken Apostelen taler, seer jeg til Nød, det vil sige: jeg
skimter ham som igjennem en Taage (thi jeg har ellers aldrig før
havt det Syn, at den almindelige Kirke virkelig skulde være den
Enkelte); men de smaa Enkelte, om hvilke Dr.\ M. taler, nei, om
jeg saa fik syv Par Briller paa: disse smaa gjenstridige Enkelte
kan jeg sandelig ikke see. Er her da et nyt Kjætteri i Gang? Vil
„den Enkelte“ maaskee have Kirkelæren omgjort, eller Dogmerne op-
% ---- printed p.56 (PDF 64) ----
løste, eller nogetsomhelst Punkt i den evangelisk-lutherske
Troesbekjendelse forvansket? Hvis Dogmatikeren nogensinde skulde
træffe paa en saadan Kjætter, da bør han visselig ikke undlade at
see ham ordentlig paa Fingrene og bekjæmpe ham grundig; thi en
Lærer i Kirken gjør vel i at give Agt paa sig selv og paa Læren
(1 Tim.\ 4, 16). Men hvis der nu ingen saadan Kjætter findes? Ja,
saa gjør Dogmatikeren vel i at give Agt paa sig selv, og lade „de
Enkelte“ være med Fred.

Hvad er det altsaa Apostelen Paulus, ifølge Dr.\ M.s
„Hermeneutik“, da egentlig vil lære os om hiin „store Enkelte“?
--- „At han er i en stadig Fremskriden henimod Troens og
Erkjendelsens Eenhed.“ --- Hvilken Erkjendelse? Thi det kan dog
vel neppe være den existentielle Troeserkjendelse; medmindre man
vil paastaae, at Christenheden i det nittende Aarhundrede har en
langt inderligere og sandere Troeserkjendelse, end paa Apostlenes
Tid. Det være langt fra mig at ville tillægge Dr.\ M. en saa
urimelig Mening, især da han jo (navnlig i §§ 185---189) selv paa
det Bestemteste hævder Inspirationen for hele den første Periode.
Kan det nu ikke være den subjective Troeserkjendelse, hvorom her
tales, nu, saa maa det jo være den objective og \textit{in specie}
den speculative. Skulde det være Meningen, da kan man uden
Betænkning give Forfatteren Ret i, at „den christelige
Speculation“ er i det nittende Aarhundrede langt videre, end i det
første, og at vi forsaavidt heri see „en stadig Fremskriden“; men
hvad følger saa heraf?
% ---- printed p.57 (PDF 65) ----
At „de Enkelte“ ikke maae være for absolute i Viden. Og hvorfor
ikke? Fordi De, der for tidligt blive absolute i Viden, ikke give
sig Stunder til at see, at den speculative Viden skal
understøttes af Troen. Men hvorfor maae saa de Enkelte ikke være
absolute i Troen? Nei, fordi De, der for tidligt blive absolute i
Troen, naturligviis forglemme, at Troen skal understøttes af den
speculative Viden. Den Enkelte, det vil sige: den lille Enkelte,
tør altsaa ikke længere erklære sig ubetinget enten for Troen
eller for Viden, tør ikke længere i nogen Henseende tale absolut
om det Absolute, efterdi „hiin store Enkelte“, ifølge Apostelens,
eller rettere ifølge Dr.\ Martensens Vidnesbyrd, endnu ikke selv
har naaet det Absolute, men „befinder sig i en fortsat
Fremskriden henimod Troens og Erkjendelsens Eenhed“. Altsaa:

\begin{center}
\emph{a) „Den christelige Dogmatik“ famler efter Problemet, men
finder det ikke.}
\end{center}

At Spørgsmaalet om Troens Forhold til Erkjendelsen staaer i nøie
Forbindelse med Spørgsmaalet om Dogmatikens Forhold til
Philosophien, navnlig til Religionsphilosophien, har Dr.\ M. meget
godt indseet; som han jo overhovedet meget godt har seet, hvilke
Punkter man i en saa vanskelig Undersøgelse især kunde fremhæve.
Strax i Indledningens § 2 erklærer han sig da ogsaa paa det
Bestemteste for en Adskillelse af \emph{Dogmatiken og
Philosophien.}

„Dogmatiken“ --- hedder det udtrykkelig --- „er ikke blot en
Videnskab om Troen, men ogsaa en Erkjendelse
% ---- printed p.58 (PDF 66) — NB: folio MISPRINTED "53" in the original ----
i Troen og udaf Troen. Den er ikke blot en Erkjendelse af, hvad
der har været, eller er Sandhed for Andre, uden at være det for
Fremstilleren selv; men heller ikke er den en philosophisk
Erkjendelse af den christelige Sandhed, hvilken Fremstilleren
havde vundet ved at indtage et Standpunkt udenfor Troen og udenfor
Kirken. Thi selv naar vi antage --- hvad vi dog ingenlunde
indrømme, --- at en videnskabelig Indsigt i den christelige
Sandhed er mulig uden christelig Tro, saa vilde en saadan
Philosopheren over Christendommen, selv om dens Resultater vare
Kirken nok saa gunstige, dog ikke kunne kaldes Dogmatik.
Dogmatiken staaer i Christendommen, og Dogmatikeren er kun sin
Videnskabs Organ, idet han tillige er sin Kirkes Organ, hvilket
ikke gjælder om den blotte Philosoph, der kun vil tjene den rene
Videnskab alene.“

Da jeg læste dette Sted første Gang, forekom det mig, at det baade
var fremsat med Klarhed, Bestemthed og Skjønhed, og at jeg
virkelig ogsaa forstod, hvad det var, der blev sagt; men --- det
være nu, som det vil --- da jeg læste det anden Gang, forstod jeg
det dog alligevel ikke. Opgavens Vanskelighed kan udtrykkes
saaledes: der leveres (som fordret% <-- print "forbret" (d/b glyph); corrected
{} er i Indledningens § 1) en \emph{videnskabelig Fremstilling og
Begrundelse af de christelige Troeslærdomme i deres Sammenhæng};
hvilken Betydning bør der nu tillægges en saadan Fremstilling?
Svar: Hvis Forfatteren er en Philosoph, der kun vil tjene den rene
Videnskab, saa har
% ---- printed p.59 (PDF 67) ----
Fremstillingen, selv om dens Resultater ere Kirken nok saa
gunstige, dog derfor endnu ikke Rang med Dogmatik; men, hvis
Forfatteren derimod vil vedføie den mundtlige eller skriftlige
Erklæring, at han sandelig har udarbeidet sin „videnskabelige
Fremstilling“ nærmest, ja allernærmest for Kirkens Skyld, ja at
han „kun er sin Videnskabs Organ, forsaavidt han er sin Kirkes
Organ“: ja saa bliver det ved Hjælp af Erklæringen og med Kirkens
Bistand og med Guds Hjælp jo nok til Dogmatik, thi Dogmatiken
staaer i Kirken, og --- Kirken staaer i Dogmatikeren. Men hvad om
„den videnskabelige Fremstilling“ nu var af en Anonym; hvad om den
anonyme Forfatter havde vidst at skjule sig saa snildt under
Personificationen af „hiin store Enkelte“, at hverken „den
kritiske Indledning“ eller „Hermeneutiken“ var istand til med
Bestemthed at afgjøre, om han nu virkelig befandt sig indenfor
Kirken eller udenfor Kirken: hvad saa? Dog dette er kun en
Hypothese; Stedet selv er hypothetisk; det vil jo \emph{ingenlunde
indrømme}, at en videnskabelig Indsigt i den christelige Sandhed
er mulig uden christelig Tro. Men hvorledes forstaae vi saa den
Sammenhæng, at hvad der ingenlunde kan indrømmes S. 4, dog nok
alligevel kan indrømmes S. 14?

„Den religiøse Anskuelse“ --- hedder det i Anm.\ S. 14 --- kan
ogsaa besiddes paa anden Haand, d.\ v.\ s.\ paa æsthetisk eller
philosophisk Maade. Og just fordi den religiøse Anskuelse bevæger
sig i Objectivitetens, i Tankens og Phantasiens Element, kan den
løsrives fra sin Livsrod i Gemyttet, og uafhængig af den
personlige Tro
% ---- printed p.60 (PDF 68) ----
besiddes paa en blot æsthetisk og philosophisk Maade. Saaledes ere
der Philosopher, Digtere, Malere, Billedhuggere, som med stor
plastisk Kraft have fremstillet de christelige Forestillinger,
uden dog selv at have besiddet disse Forestillinger paa religiøs
Maade, idet de kun have forholdt sig dertil igjennem Tankens og
Phantasiens Medium.“

Det er ret Skade, at disse to vigtige Steder (S. 4 og S. 14) holde
sig saa vidt adskilte i Bogen; havde de staaet hinanden nærmere,
vilde Forfatteren uden Tvivl være bleven opmærksom paa den
tilsyneladende Modsigelse, og for Læserens Skyld have tilføiet en
lille oplysende Mellemsætning. Som de nu staae, ligne de to
Paarørende, der have skilt sig i indbyrdes Uenighed og vende
hinanden Ryggen. --- Har man nemlig først indrømmet, at der ere
Philosopher, som uden selv at være religiøse (ɔ: uden sand
christelig Tro) have fremstillet de christelige Forestillinger med
stor plastisk Kraft (altsaa med philosophisk, det vil sige:
videnskabelig Indsigt i disse Forestillingers Sandhed); nu, saa
kunde man jo dog, om ikke for Andet, saa dog for Conseqvensens
Skyld, ligesaa gjerne indrømme, at en videnskabelig Indsigt i den
christelige Sandhed er mulig uden christelig Tro; men nu er det jo
just Det, man „ingenlunde vil indrømme.“

At Dr.\ M. paa det Bestemteste vil have Dogmatiken adskilt fra
Religionsphilosophien, er tydeligt nok, vel ikke just af
Udførelsen, men dog af Indledningen. End ikke den christelige
Philosophie tør her falde sammen
% ---- printed p.61 (PDF 69) ----
med den dogmatiske Speculation. Det hedder (S. 81 § 35) blandt
Andet saaledes:

„Men efterat Christendommen selv har født en ny Erkjendelse af
sig, en Theologie, opstaaer det Spørgsmaal, om der ved Siden af
Theologien ogsaa bliver Plads for en christelig Philosophie, og
hvorledes disse da forholde sig til hinanden. Vi forudsætte her,
at der gives en christelig Philosophie, vi forudsætte fremdeles,
at den er bunden til de samme Erkjendelsens Grundbetingelser som
Theologien, altsaa maa gaae ud fra \textit{credo ut intelligam};
men vi sætte Forskjellen deri, at Philosophien ogsaa som christelig
er Verdens-Viisdom, d.\ v.\ s.\ Universumserkjendelse, medens
Theologien er Gudserkjendelse, en Forskjel, der vistnok er
relativ, men dog en Forskjel.“

Ja vistnok er denne Forskjel imellem den christelige Philosophie
og den dogmatiske Speculation en relativ Forskjel; den er
ingenlunde absolut, men relativ, kun relativ: det er et sandt Ord,
det er godt sagt; men lad os nu undersøge, om her ogsaa virkelig
er nogen Forskjel, eller om den relative Forskjel ikke tilsidst
kryber ind og bliver mindre og mindre, om den da ikke tilsidst
bliver saa lille og fiin, at Ingen kan vise den frem, og Ingen
faae den at see, fordi den paa Grund af sin uendelige Relativitet
tilsidst befindes at være aldeles bortkommen og forsvunden.

Forskjellen bestemmes saaledes (S. 82): „Philosophien søger at
erkjende den guddommelige Verdenslov, der
% ---- printed p.62 (PDF 70) ----
fuldbyrder sig igjennem hele Tilværelsen, gjennem Naturens og
Aandeverdenens forskjellige Kredse, og den søger at erkjende
Christendommen som Verdenslovens Fylde. Philosophien fordyber sig
altsaa i det Mangfoldige og fører dette tilbage til det Ene, til
Guds Rige som det altopklarende Midtpunkt; Theologien, Dogmatiken
derimod tager fra først af sit Standpunkt i Centrum, fordyber sig
udelukkende i det Ene, i Guds Rige som saadant.“ --- Men, idet vi
nu ogsaa fordybe os i, hvad der her er sagt om „det Ene,“ om „Guds
Rige som saadant“, bør vi da heller ikke undlade at fordybe os i
alt Det, Dr.\ M. paa saa mange andre Steder har sagt om „det
kosmiske Under“, om „Skabelse og Kosmogonie“, om „Mennesket og
Englene“, om „den syndige Naturbestemthed og den syndige
Historie“, om „de dæmoniske Magter og Djævelen“, om „Døden og
Alskabningens Forkrænkelighed“, om „det frie Verdensløb og Guds
mangfoldige Viisdom“, om „Faderens Logos“, som ikke blot er „Gud
Faders Hjerte“, men tillige det evige „Verdenshjerte“ o.\ s.\ v.\
o.\ s.\ v.; thi alt Dette viser tilbage til hvad der udtrykkeligt
er os sagt (S. 13) om Religion og Aabenbaring, at nemlig „Ideen om
Gud vinder Skikkelse i en omfattende \emph{Anskuelse} af Verden og
Menneskelivet i dets Forhold til Gud, en Anskuelse af Himmel og
Jord, Natur og Historie, Himmel og Helvede.“

Er nu den religiøse Anskuelse virkelig saa omfattende, og skal alt
Dette gjennemtænkes, udforskes og fremstilles med videnskabelig,
ja med speculativ-videnskabelig Grundighed; da

% ---- printed p.63 (PDF 71) ----
synes det unegtelig at fremgaae heraf, at en saa omfattende
dogmatisk Gudserkjendelse umulig kan adskilles fra en ligesaa
omfattende philosophisk Verdenserkjendelse, ɔ: en
dialektisk-speculativ Universumserkjendelse. Men hvilken Forskjel
er der saa imellem Dogmatiken og Christendomsphilosophien? Svar:
ingen, aldeles ingen! Philosophen maa tage sit Standpunkt i
„Centrum“ ligesaavel som Dogmatikeren; thi hvorledes skulde han
vel erkjende Christendommen som „Verdenslovens Fylde“, uden at
erkjende den som „det Centrale“; men hvorledes skulde han
erkjende det Centrale, uden at tage sit Standpunkt i Centrum? Det
er forgjæves, at Dogmatikeren beraaber sig paa Troen og puster sig
op i Kirken, naar han% <-- print "har" (n/r glyph); corrected
{} siger speculativt: ποῦ στῶ, thi Philosophen siger nu ogsaa
speculativt: πᾷ στῶ; og see, saa staae de jo begge lige midt i
Speculationen. Kan Dogmatikeren ɔ: den speculerende Dogmatiker
ikke undvære „de nye Verdenskategorier“, saa kan
Religionsphilosophen da heller ikke forsømme „de gamle
Aabenbaringskategorier“; thi den speculative Productivitet er just
betinget ved begges gjensidige Udvikling. „Den christelige
Speculation“, der jo netop gaaer ud paa, at forklare Universet i
Christendommen og omvendt Christendommen i Universet, kan umulig
adskille disse to Functioner, uden at opløse og forqvakle sig
selv.

„Men“ --- siger man --- „der er dog alligevel en Forskjel, nemlig
i Stoffets Behandling er der en Forskjel, en ikke ubetydelig
Forskjel; thi, medens Dogmatikeren, der ikke blot har πᾷ στῶ i
Speculationen, men desfor-
% ---- printed p.64 (PDF 72) ----
uden et ganske andet, høiere og langt tilforladeligere πᾷ στῶ i
Kirken, gaaer i Detail'en og tager det nøiagtig med de enkelte
positive Læresætninger; pleier Philosophen derimod kun at behandle
Sagen en gros og forbigaae mange kirkelige Punkter, blot for at
faae Christendomsphilosophien med op som Led i sin philosophiske
Encyklopædie.“ --- Men i Forhold til Principet er dette jo en
aldeles intetsigende Forskjel. Thi, hvis den speculerende
Dogmatiker optager mere Stof end Speculationen tillader,
forqvakler han jo Videnskaben; og omvendt, hvis Philosophen
forsømmer at tilegne sig alt det Stof, Speculationen kræver,
forqvakler han jo ligeledes Videnskaben: hvor er saa Forskjellen?

„Ja - men - jo“ --- siger man --- „her er dog alligevel en
Forskjel, en vis Forskjel, jo vist, er her en Forskjel, det kan
ikke negtes.“ --- Hvori bestaaer da Forskjellen? --- „Hvori den
bestaaer? Ja det kan naturligviis ikke saadan siges med
afgjørende Bestemthed, da den sunde Betragtning jo strax viser, at
den Grændse, vi her have for os, maa være en flydende Grændse, og
forsaavidt ogsaa en vordende Grændse. Den sunde Betragtning
(jfr.\ S.\,80) vil altid erkjende, at den speculative Begriben
selv er et saare bevægeligt og dialektisk Begreb, der ikke lader
sig affærdige med et tørt Ja eller Nei, ikke lader sig affærdige
med den Paastand, at den enten fuldstændigt maa være eller ikke
være; thi den er kun som \emph{vordende}. Saameget kan dog
imidlertid siges med Bestemthed, at Dogmatikeren staaer i Kirken,
og at han
% ---- printed p.65 (PDF 73) ----
kun er sin Videnskabs Organ, idet han tillige er sin Kirkes Organ,
hvilket ikke gjælder om den blotte Philosoph, der kun vil tjene
den rene Videnskab alene. Staaer dette nu fast, at Dogmatikeren
staaer i Kirken, og elsker Troen som han elsker Kirken, saa bør
det ogsaa siges med Bestemthed, at den intellectuelle Kjærlighed
til den christelige Sandhed, der navnlig maa findes hos
\emph{Lærerstanden i Kirken}, er langt høiere, langt personligere
og langt mere pathologisk, end den, der pleier at findes hos den
blotte Philosoph.“ --- Men, naar det nu er vitterligt, at denne
Adskillelse beroer paa en ulyksalig Forvexling% <-- print "Forverling" (r/x glyph); corrected
{} af det Personlige og det Geniale; naar det nu er vitterligt, at
En kan være en troende Christen og arbeide kraftig i Troen, uden
at have den intellectuelle, d.\ v.\ s.\ den theoretiske
Kjærlighed, hvorom her tales, medens en Anden kan være aldeles
behersket og betagen af denne objective Kjærlighed, uden i samme
Grad at arbeide for den subjective Tilegnelse i Troen; naar det er
vitterligt, --- hvad Forfatteren af „den christelige Dogmatik“
tilvisse veed saa godt, som Nogen --- at den intellectuelle og
speculative Kjærlighed, hvormed Tænkeren fordyber sig i
Aabenbaringens Mysterier, ingenlunde er nogen egentlig religiøs
Fordring („Troen er jo det ene Fornødne for os Alle“), men ene og
alene Talentets Drift og Trang til Beskuelse af Ideen, til
objectiv Glæde i guddommelige Syner: hvor er saa Forskjellen?

Vistnok kunne vi ikke negte Dr.\ M. den Anerkjendelse, at han, som
dogmatisk Genius, har været betænkt
% ---- printed p.66 (PDF 74) ----
paa at gjøre en bestemt Adskillelse imellem Dogmatiken og
Religionsphilosophien (en Adskillelse, der som bekjendt er af en
overordentlig Vigtighed); men saa meget mere maae vi beklage, at
det endnu ikke er lykkedes dette Talent at finde det Punkt,
hvorfra Adskillelsen gaaer ud. De Exempler, han leilighedsviis
anfører, ere saa uheldige, som tænkes kan. For at vindicere
Dogmatiken sit selvstændige Princip, viser han hen til
Schleiermacher (S. 7); men hvad kan det hjælpe, naar
Schleiermachers Standpunkt (S. 11 § 7) er forfeilet, og hans
dogmatiske Opfattelse af Religionen et vistnok meget genialt og
dybsindigt, men dog mislykket Forsøg, der hverken stemmer eller
kan stemme med den orthodoxe Kirkelære? For at udtrykke det
speculative Objectivitetsforhold til Christendommen, henviser
Forfatteren (S. 76) til „den speculative Mystik og Theosophie
(Jacob Böhme), som lærer, at der i Troen selv er en Skuen“; men
--- ret som om det dog alligevel var forudbestemt, at Læseren
hverken skulde have Gavn eller Glæde af denne Skuen --- kalder han
igjen Henviisningen tilbage, og gjør os opmærksomme paa, at ogsaa
her var en „Misviisning, der viste bort fra Historien.“ Hvor alt
dette dog er uheldigt, ja ubehageligt! Skal Dogmatiken med
Bestemthed adskilles fra Religionsphilosophien, da maa Forskjellen
absolut kunne angives saaledes, at disse Videnskaber sees enten at
have en ganske forskjellig Opgave, eller, hvis Opgaven paa et vist
Punkt befindes at være den samme, da at løse den paa en saa
forskjellig Maade, at Methoderne idet-
% ---- printed p.67 (PDF 75) ----
mindste divergere, som jo Vinkelens Sider divergere, endskjøndt de
begynde fra det samme Toppunkt. Men --- medmindre jeg
ulykkeligviis har seet mig blind paa al den speculative Glands,
der er udgydt over den martensenske Indledning --- den paatænkte
Adskillelse ender i et reent Nul; thi Standpunktet, det vil sige:
det videnskabelige Standpunkt (hvad Dogmatikeren for Resten
bestiller i Kirken, vedkommer ikke denne Sag) bliver det samme,
Opgaven den samme, Løsningen den samme.

Og Dr.\ M. vil føre Ordet for „den sammenhængende Erkjendelse“!
Det kan nu være meget godt med den sammenhængende Erkjendelse og
„den flydende Grændse“, naar Alt forøvrigt vil forme sig, naar
Dialektiken arbeider flink og allevegne virker den behørige
Sondring; men naar Forskjellene flyde sammen i det Flydende,
medens Sammenhængen fortætter sig i et mystisk Halvmørke, der
skjuler Problemet: da bør man visselig ikke forsmaae en god
„Aphorisme“, en „Strøtanke“ eller et „Glimt“, især, hvis det er et
Lynglimt.

\begin{center}
\emph{b) „Den christelige Dogmatik“ vil tilegne sig Speculationen,
men uden at fatte Speculationens Problem.}
\end{center}

Det kan visselig ikke miskjendes, at Dr.\ M. jo paa mange Maader
og under mange forskjellige Vendinger stræber at indprænte
Dogmatikens Læsere Uundværligheden af en objectiv Viden. I denne
Charakteer optræder han saavel paa Ideens som paa Kirkens Vegne
som Talsmand for den høiere Videnskab og navnlig for
% ---- printed p.68 (PDF 76) ----
Speculationen. Det er et hæderligt Hverv, en smuk Stilling, især i
et lille Land, hvor deres Tal, der føle Trang og Kald til grundig
at give sig af med speculativ Philosophie, ordentligviis ikke kan
være meget stort. Imod denne Optræden er her da i og for sig,
saavidt jeg skjønner, Intet at indvende; og jeg for mit
Vedkommende vilde gjerne efter bedste Evne staae Dr.\ M. bi i hans
agtværdige Stræben. Men hvad der i dette Tilfælde ængster mig
noget, er det skarpe Sideblik, hvormed han betragter visse
„Enkelte“, som, efter først at have forsøgt sig lidt i
Speculationen, og derved „altfor tidligt“ at være blevne „absolute
i Viden“, nu dernæst i den senere Tid have kastet Speculationen
reent bort, og „altfor tidligt“ ere blevne „absolute i Troen“. Det
er ret en fatal Omstændighed med slige forblommede Yttringer, der
gjør En saa uvis, saa forlegen, naar man just har den bedste
Villie til at tjene den bedste Sag. Man har en Angest for at
tilbyde sin Tjeneste, naar man ikke ret veed, hvorledes den vil
blive modtaget: maaskee man bliver afviist som --- ubrugbar,
maaskee man allerede er afviist; maaskee „den christelige
Viisdom“\footnote{Thi „den christelige Viisdom“ --- ja det maa jeg
tilstaae --- „den christelige Viisdom“ veed at pointere.} allerede
har casseret En som Efterligner, som --- \textit{miserabile dictu}
--- „trivial Efterligner“, og ladet En afgaae ved sidste
Reduction; maaskee man saa engang hændelsesviis faaer det at vide,
at der er
% ---- printed p.69 (PDF 77) ----
En, der er afgaaet fra sin Videnskab, ligesom Manden, der læste i
Avisen, hændelsesviis fik at vide, at det var ham, der var afgaaet
fra sit Embede. Det er en fatal Omstændighed; thi hvordan man end
snoer og dreier det, noget Mystisk blander sig dog heri. Jeg vil
da tage Sagen saa ligefrem og eenfoldig, som muligt: er det
virkelig mig, der er bleven dømt, reduceret og casseret, afsat fra
„den sammenhængende Tænkning“ og entlediget fra „den objective
Viden“: ja, saa faaer jeg at finde mig deri, og det skal heller
ikke komme derpaa an, naar blot min Dom maa blive mig offentlig
forkyndt, at jeg kan faae min Skjæbne at vide. Men ville de
forblommede Allusioner kanskee ikke slaae mig reent ihjel; maaskee
jeg da foreløbig kan redde mig ud af Forlegenheden ved at tilføie
en Bemærkning, min philosophiske Stræben betræffende. I de ti Aar,
jeg nu saadan \textit{con amore et ex officio} har beskjæftiget
mig med Philosophien, den nyere og den ældre, har jeg med særdeles
Interesse fulgt Speculationen ad dens forskjellige Veie og derved
havt Leilighed til at see, hvorvidt Metaphysikens hidtil opdagede
Enemærker vel omtrent kunne strække sig. Har jeg end ikke selv
udfundet Noget, der er værd at omtale, saa har jeg dog hos mine
Læremestere faaet et stort Indtryk af, hvad „sammenhængende
Erkjendelse“ har at betyde, og --- hvad jeg betragter som det
glædeligste Udbytte af mine Studier --- lært at beundre den
systematiske Tænknings sande Heroer og at bøie mig for Geniets
Overlegenhed. Skjøndt jeg nu vistnok ikke er noget afgjort
philosophisk Talent
% ---- printed p.70 (PDF 78) ----
(endnu mindre en Genius), tiltroer jeg mig dog saameget Hoved, at
jeg endda nogenledes kan gjøre Forskjel paa Speculation og
Phantasie, og ret godt følge med, især naar en bedre Begavet vil
føre an. Saa lader jeg da igjen Dr.\ M. føre an --- med det
speculative Thema, og følger ham, som jeg kan, for dog at
accompagnere med lidt Dialektik.

Mange og mærkværdige ere de Steder, der i „den christelige
Dogmatik“ tale til Speculationens Ophøielse. „Thi“ --- hedder det
(S. 78) --- „skjøndt vi fuldkomment indrømme, hvad Melanchthon
\ldots{} siger mod en ørkesløs Speculation, der var løsreven fra
Livet, saa er dog Erkjendelsen af Guds Hemmeligheder i sig selv et
Gode, og Erkjendelsen af Guds Herlighed er i sig selv
opbyggelig.“ Endvidere (S. 79): „Den grundige Explication vil ikke
kunne Andet end udvikle saadanne Tanke-Modsætninger, saadanne
Antinomier, der kræve en \emph{Mediation i Begrebet}; thi, som det
hedder i Sirachs Viisdom (33,17): „Den Høiestes Gjerninger ere
stedse To, een imod den anden“; og det Speculative beroer netop
paa at fatte Modsætningerne i Ideens Eenhed. En speculativ Skuen
maa stedse forudsættes, naar Fremstillingen ikke skal tabe sig i
udvortes Forstandighed, eller kun indskrænke sig til at opfatte de
christelige Forestillinger i deres praktiske, blot anvendelige
Betydning“ --- Derfor skulle ogsaa alle De have Utak, der (S. 85)
„skjelne mellem Theologie og Philosophie, som mellem rene og urene
Spiser, der sige om Philosophien: Smag

% [text to be added: pp. 71--132]

\end{document}
